\documentclass[12pt,a4paper]{article}

\usepackage[a4paper, top=2.5cm, bottom=2.5cm, left=3cm, right=2.5cm]{geometry}
\usepackage{amsmath, amssymb, mathtools}

\title{Slovní úlohy}
\author{Ondřej Škubala}
\date{\today}

\begin{document}

\maketitle


% ---------------------------------
% ------ Ú L O H A  -  č. 1 -------
% ---------------------------------
\newpage
\section*{Úloha č. 1}

\textbf{Zadání:}  
Žáci jedné třídy chtějí paní učitelce věnovat lístek do divadla. Jestliže každý z nich přispěje 12 korunami, k zakoupení lístku jim bude chybět 34 korun. Přispěje-li každý žák 15 korunami, zbude jim 50 Kč. Nakonec se žáci dohodli, že každý přinese 14 korun.

Vypočtěte, kolik korun třídě zbude po zakoupení lístku.  
V záznamovém archu uveďte celý postup řešení.

\bigskip

\textbf{Řešení:}

Označme si:
\begin{itemize}
    \item \(x\) \dots počet žáků,
    \item \(y\) \dots cena lístku.
\end{itemize}

Pokud každý žák přispěje 12 Kč, vyberou dohromady
\[
12x.
\]
To ale nestačí, a bude jim chybět 34 Kč. Platí tedy:
\[
12x + 34 = y.
\]

Pokud každý žák přispěje 15 Kč, vyberou dohromady
\[
15x.
\]
Tentokrát už budou mít více, konkrétně o 50 Kč. Platí tedy:
\[
15x - 50 = y.
\]

Obě rovnice vyjadřují stejnou cenu lístku, proto je položíme sobě rovny:
\[
12x + 34 = 15x - 50.
\]

Rovnici upravíme:
\[
34 + 50 = 15x - 12x,
\]
\[
84 = 3x,
\]
\[
x = 28.
\]

Ve třídě je tedy \(28\) žáků.

Nyní určíme cenu lístku:
\[
y = 12 \cdot 28 + 34 = 336 + 34 = 370.
\]

Cena lístku je tedy \(370\) Kč.

Nakonec se žáci dohodli, že každý přinese 14 Kč. Vyberou proto celkem
\[
14 \cdot 28 = 392.
\]

Po zakoupení lístku jim zbude
\[
392 - 370 = 22.
\]

\bigskip

\textbf{Odpověď:}  
Po zakoupení lístku třídě zbude 22 Kč.


% ---------------------------------
% ------ Ú L O H A  -  č. 2 -------
% ---------------------------------
\newpage
\section*{Úloha č. 2}

\textbf{Zadání:}  
Pětimístné přirozené číslo je sestaveno z pěti různých číslic. Uprostřed je vždy číslice 6. Všechny číslice jsou seřazeny sestupně, tedy od největší po nejmenší.

(Daným podmínkám vyhovují např. čísla 97650 a 87631.)

Kolik různých čísel je možné uvedeným způsobem sestavit?

\bigskip

\textbf{Řešení:}

Hledáme pětimístná čísla, která mají tvar
\[
\overline{ab6de},
\]
kde číslice jsou seřazeny sestupně. To znamená, že musí platit
\[
a>b>6>d>e.
\]

Z toho plyne:
\begin{itemize}
    \item vlevo od šestky musí být dvě různé číslice větší než 6,
    \item vpravo od šestky musí být dvě různé číslice menší než 6.
\end{itemize}

\medskip

Číslice větší než 6 jsou:
\[
7,8,9.
\]
Z těchto tří číslic vybíráme dvě:
\[
\binom{3}{2}=3.
\]

Číslice menší než 6 jsou:
\[
0,1,2,3,4,5.
\]
Z těchto šesti číslic vybíráme dvě:
\[
\binom{6}{2}=15.
\]

Jakmile číslice vybereme, jejich pořadí už je jednoznačně určeno, protože celé číslo musí být zapsáno sestupně.

Celkový počet různých čísel je tedy
\[
3 \cdot 15 = 45.
\]

\bigskip

\textbf{Odpověď:}  
Uvedeným způsobem je možné sestavit 45 různých čísel.

% ---------------------------------
% ------ Ú L O H A  -  č. 3 -------
% ---------------------------------
\newpage
\section*{Úloha č. 3}

\textbf{Zadání:}  
V letadle na cestě do Asie letělo o 40 cizinců více než Čechů. Polovina cizinců požadovala vegetariánskou stravu, z českých pasažérů měla stejné přání pouze desetina. Vegetariánská strava se tak připravovala pro třetinu všech pasažérů.

Určete celkový počet pasažérů požadujících vegetariánskou stravu.

\bigskip

\textbf{Řešení:}

Označme si:
\begin{itemize}
    \item \(x\) \dots počet Čechů,
    \item \(x+40\) \dots počet cizinců.
\end{itemize}

Podle zadání požadovala vegetariánskou stravu:
\begin{itemize}
    \item polovina cizinců, tedy
    \[
    \frac{1}{2}(x+40),
    \]
    \item desetina Čechů, tedy
    \[
    \frac{1}{10}x.
    \]
\end{itemize}

Celkový počet pasažérů byl
\[
x+(x+40)=2x+40.
\]

Vegetariánskou stravu požadovala třetina všech pasažérů, tedy
\[
\frac{1}{3}(2x+40).
\]

Sestavíme rovnici:
\[
\frac{1}{2}(x+40)+\frac{1}{10}x=\frac{1}{3}(2x+40).
\]

Nyní rovnici upravíme:
\[
\frac{x}{2}+20+\frac{x}{10}=\frac{2x+40}{3}.
\]

Levou stranu sečteme:
\[
\frac{5x}{10}+\frac{x}{10}+20=\frac{6x}{10}+20=\frac{3x}{5}+20.
\]

Dostáváme:
\[
\frac{3x}{5}+20=\frac{2x+40}{3}.
\]

Celou rovnici vynásobíme 15:
\[
15\left(\frac{3x}{5}+20\right)=15\left(\frac{2x+40}{3}\right).
\]

\[
9x+300=10x+200.
\]

\[
300=x+200,
\]

\[
x=100.
\]

Počet Čechů je tedy
\[
100.
\]

Počet cizinců je
\[
100+40=140.
\]

Nyní určíme počet pasažérů požadujících vegetariánskou stravu:
\[
\frac{1}{2}\cdot 140+\frac{1}{10}\cdot 100=70+10=80.
\]

\bigskip

\textbf{Odpověď:}  
Vegetariánskou stravu požadovalo celkem 80 pasažérů.


% ---------------------------------
% ------ Ú L O H A  -  č. 4 -------
% ---------------------------------
\newpage
\section*{Úloha č. 4.1}

\textbf{Zadání:}  
Hráč přichází na stanoviště s určitým počtem žetonů. Na stanovišti musí utratit alespoň jeden žeton.

Kolik žetonů hráč utratí, tolikrát se při odchodu ze stanoviště zvětší počet jeho zbývajících žetonů. (Např. přichází-li hráč na stanoviště s 10 žetony a utratí 4 žetony, stanoviště opustí s 24 žetony.)

Žetony nelze dělit.

Aleš přichází na stanoviště se 45 žetony.

Určete počet žetonů, které musí Aleš utratit, aby stanoviště opouštěl nejméně s 500 žetony. Najděte všechna řešení.

\bigskip

\textbf{Řešení:}

Označme si
\[
x
\]
počet žetonů, které Aleš utratí.

Po utracení \(x\) žetonů mu zbude
\[
45-x
\]
žetonů.

Podle zadání se při odchodu tento počet zvětší \(x\)-krát. Aleš tedy odejde se
\[
x(45-x)
\]
žetony.

Požadujeme, aby odcházel nejméně s 500 žetony. Sestavíme nerovnici:
\[
x(45-x)\geq 500.
\]

Roznásobíme:
\[
45x-x^2\geq 500.
\]

Převedeme vše na jednu stranu:
\[
-x^2+45x-500\geq 0.
\]

Vynásobíme nerovnici číslem \(-1\), přičemž se obrátí znaménko:
\[
x^2-45x+500\leq 0.
\]

Nyní vyřešíme kvadratickou rovnici
\[
x^2-45x+500=0.
\]

Diskriminant je
\[
D = (-45)^2-4\cdot 1\cdot 500 = 2025-2000=25.
\]

Kořeny rovnice jsou:
\[
x_{1,2}=\frac{45\pm \sqrt{25}}{2}=\frac{45\pm 5}{2}.
\]

Dostáváme:
\[
x_1=\frac{40}{2}=20,\qquad x_2=\frac{50}{2}=25.
\]

Protože výraz
\[
x^2-45x+500
\]
má být menší nebo roven nule, platí:
\[
20\leq x\leq 25.
\]

Žetony nelze dělit, proto \(x\) musí být celé číslo. Vyhovují tedy hodnoty:
\[
x\in\{20,21,22,23,24,25\}.
\]

\bigskip

\textbf{Odpověď:}  
Aleš musí utratit 20, 21, 22, 23, 24 nebo 25 žetonů.

\section*{Úloha č. 4.2}

\textbf{Zadání:}  
Určete největší možný počet žetonů, který si Aleš z tohoto stanoviště může odnést.

\bigskip

\textbf{Řešení:}

Stejně jako v předchozí části označme
\[
x
\]
počet utracených žetonů.

Počet žetonů, se kterými Aleš odejde, je
\[
x(45-x).
\]

Po úpravě dostaneme:
\[
x(45-x)=45x-x^2.
\]

Jde o kvadratickou funkci otevřenou dolů, takže její maximum leží ve vrcholu paraboly.

Souřadnice vrcholu určíme:
\[
x_v=\frac{-b}{2a}=\frac{-45}{2(-1)}=\frac{45}{2}=22{,}5.
\]

Protože \(x\) musí být celé číslo, vyzkoušíme dvě nejbližší hodnoty:
\[
x=22 \quad \text{a} \quad x=23.
\]

Pro \(x=22\):
\[
22(45-22)=22\cdot 23=506.
\]

Pro \(x=23\):
\[
23(45-23)=23\cdot 22=506.
\]

V obou případech dostáváme stejný maximální počet žetonů.

\bigskip

\textbf{Odpověď:}  
Největší možný počet žetonů, který si Aleš může odnést, je 506.


% ---------------------------------
% ------ Ú L O H A  -  č. 5 -------
% ---------------------------------
\newpage
\section*{Úloha č. 5}

\textbf{Zadání:}  
Dlažba kolem stožáru na vlajku (černý otvor) vytváří pravidelné tmavé a světlé prstence.  
(Všechny dlažební kostky jsou shodné pravidelné šestiboké hranoly.)

Kolik prstenců je vytvořeno z \(1\,260\) dlaždic?

\bigskip

\textbf{Řešení:}

Z obrázku vidíme, že:
\begin{itemize}
    \item 1 prstenec je vytvořen z \(6\) dlaždic,
    \item 2 prstence jsou vytvořeny z \(6+12\) dlaždic,
    \item 3 prstence jsou vytvořeny z \(6+12+18\) dlaždic.
\end{itemize}

Počet dlaždic v jednotlivých prstencích tedy tvoří posloupnost:
\[
6,\ 12,\ 18,\ 24,\dots
\]

Jde o aritmetickou posloupnost s prvním členem
\[
a_1=6
\]
a diferencí
\[
d=6.
\]

Počet dlaždic v \(n\)-tém prstenci je:
\[
a_n=6n.
\]

Celkový počet dlaždic v \(n\) prstencích je součet prvních \(n\) členů této posloupnosti:
\[
S_n=\frac{n}{2}(a_1+a_n).
\]

Dosadíme:
\[
S_n=\frac{n}{2}(6+6n).
\]

Vytkneme šestku:
\[
S_n=\frac{n}{2}\cdot 6(n+1)=3n(n+1).
\]

Podle zadání je celkem \(1\,260\) dlaždic, tedy:
\[
3n(n+1)=1260.
\]

Vydělíme třemi:
\[
n(n+1)=420.
\]

Dostáváme rovnici:
\[
n^2+n-420=0.
\]

Rozložíme na součin:
\[
(n-20)(n+21)=0.
\]

Odtud:
\[
n=20 \quad \text{nebo} \quad n=-21.
\]

Záporné řešení nemá smysl, proto platí:
\[
n=20.
\]

\bigskip

\textbf{Odpověď:}  
Z \(1\,260\) dlaždic je vytvořeno 20 prstenců.


% ---------------------------------
% ------ Ú L O H A  -  č. 6 -------
% ---------------------------------
\newpage
\section*{Úloha č. 6}

\textbf{Zadání:}  
Číslo, které se čte stejně zleva i zprava, se nazývá palindrom. Uvažujme všechny pětimístné palindromy, které mají první číslici větší než druhou.

Kolik různých palindromů je možné uvedeným způsobem sestavit?

\bigskip

\textbf{Řešení:}

Pětimístný palindrom má tvar
\[
\overline{abcba}.
\]

Podmínka ze zadání říká, že první číslice musí být větší než druhá, tedy:
\[
a>b.
\]

\medskip

Nyní určíme možnosti pro jednotlivé číslice:

\begin{itemize}
    \item první číslice \(a\) nemůže být nula, takže
    \[
    a\in\{1,2,3,4,5,6,7,8,9\},
    \]
    tedy má \(9\) možností,
    
    \item druhá číslice \(b\) musí být menší než \(a\),
    
    \item prostřední číslice \(c\) může být libovolná od 0 do 9, tedy má \(10\) možností.
\end{itemize}

\medskip

Pro každou volbu první číslice \(a\) spočítáme, kolik je možností pro \(b\):

\[
\begin{array}{c|c}
a & \text{počet možností pro } b \\
\hline
1 & 1 \\
2 & 2 \\
3 & 3 \\
4 & 4 \\
5 & 5 \\
6 & 6 \\
7 & 7 \\
8 & 8 \\
9 & 9 \\
\end{array}
\]

Celkový počet dvojic \((a,b)\) je tedy:
\[
1+2+3+4+5+6+7+8+9=45.
\]

Ke každé takové dvojici můžeme zvolit prostřední číslici \(c\) deseti způsoby, takže celkový počet palindromů je:
\[
45\cdot 10=450.
\]

\bigskip

\textbf{Odpověď:}  
Uvedeným způsobem je možné sestavit 450 různých palindromů.


% ---------------------------------
% ------ Ú L O H A  -  č. 7 -------
% ---------------------------------
\newpage
\section*{Úloha č. 7}

\textbf{Zadání:}  
Aleš zaplatil za zlevněný zájezd \(9\,000\) korun a z původní ceny zájezdu tak ušetřil čtvrtinu. Sleva se týkala jen dopravy, jejíž cena klesla na \(40\ \%\) původní ceny. Ostatní náklady zůstaly v plné výši.

Vypočtěte původní cenu dopravy \(d\).

\bigskip

\textbf{Řešení:}

Označme si:
\begin{itemize}
    \item \(P\) \dots původní cena celého zájezdu,
    \item \(d\) \dots původní cena dopravy.
\end{itemize}

Aleš zaplatil po slevě \(9\,000\) Kč a ušetřil čtvrtinu původní ceny. To znamená, že zaplatil
\[
\frac{3}{4}
\]
původní ceny zájezdu.

Platí tedy:
\[
\frac{3}{4}P=9000.
\]

Odtud:
\[
P=9000\cdot \frac{4}{3}=12000.
\]

Původní cena zájezdu byla tedy
\[
12\,000\text{ Kč}.
\]

Sleva se týkala jen dopravy. Doprava po slevě stojí \(40\%\) původní ceny, tedy
\[
0{,}4d.
\]

Ostatní náklady zůstaly stejné. Zlevnění celého zájezdu tedy vzniklo jen tím, že doprava klesla z \(d\) na \(0{,}4d\).

Úspora na dopravě je:
\[
d-0{,}4d=0{,}6d.
\]

Zároveň víme, že Aleš ušetřil čtvrtinu z původní ceny zájezdu:
\[
\frac{1}{4}\cdot 12000 = 3000.
\]

Tato úspora je právě rovna úspoře na dopravě, tedy:
\[
0{,}6d=3000.
\]

Odtud:
\[
d=\frac{3000}{0{,}6}=5000.
\]

\bigskip

\textbf{Odpověď:}  
Původní cena dopravy byla 5\,000 Kč.


% ---------------------------------
% ------ Ú L O H A  -  č. 8 -------
% ---------------------------------
\newpage
\section*{Úloha č. 8.1}

\textbf{Zadání:}  
Pětimístné číslo má ve svém zápise čtyřikrát stejnou nenulovou číslici a jednu větší číslici.

Určete, kolik čísel vyhovujících podmínkám zadání má ve svém zápise číslici 1.

\bigskip

\textbf{Řešení:}

Máme pětimístné číslo, které obsahuje:
\begin{itemize}
    \item čtyři stejné číslice,
    \item jednu větší číslici.
\end{itemize}

Podle zadání má číslo obsahovat číslici 1, takže čtyřikrát se opakuje právě číslice
\[
1.
\]

Zbývající číslice musí být větší než 1, tedy:
\[
x \in \{2,3,4,5,6,7,8,9\}.
\]

Máme tedy čísla složená ze čtyř jedniček a jedné číslice \(x\).

\medskip

Nejprve zvolíme číslici \(x\):
\[
8 \text{ možností}.
\]

Poté zvolíme pozici, na které se tato číslice nachází:
\[
5 \text{ možností}.
\]

Celkový počet čísel je tedy:
\[
8 \cdot 5 = 40.
\]

\bigskip

\textbf{Odpověď:}  
Takových čísel je 40.

\section*{Úloha č. 8.2}

\textbf{Zadání:}  
Určete počet všech čísel vyhovujících podmínkám zadání.

\bigskip

\textbf{Řešení:}

Obecně máme číslo, které obsahuje:
\begin{itemize}
    \item čtyři stejné nenulové číslice \(a\),
    \item jednu větší číslici \(b\), tedy \(b>a\).
\end{itemize}

\medskip

Nejprve volíme číslici \(a\):
\[
a \in \{1,2,3,4,5,6,7,8\}.
\]

(Pro \(a=9\) už neexistuje větší číslice.)

\medskip

Pro každé \(a\) volíme číslici \(b>a\):

\[
\begin{array}{c|c}
a & \text{počet možností pro } b \\
\hline
1 & 8 \\
2 & 7 \\
3 & 6 \\
4 & 5 \\
5 & 4 \\
6 & 3 \\
7 & 2 \\
8 & 1 \\
\end{array}
\]

Celkem:
\[
8+7+6+5+4+3+2+1=36.
\]

Pro každou dvojici \((a,b)\) můžeme umístit číslici \(b\) na libovolnou z pěti pozic:
\[
5 \text{ možností}.
\]

Celkový počet čísel:
\[
36 \cdot 5 = 180.
\]

\bigskip

\textbf{Odpověď:}  
Takových čísel je 180.


% ---------------------------------
% ------ Ú L O H A  -  č. 9 -------
% ---------------------------------
\newpage
\section*{Úloha č. 9.1}

\textbf{Řešení:}

Průměry dlaždic tvoří aritmetickou posloupnost:
\[\
51,\ 52,\ 53,\dots
\]

Platí:
\[
a_1=51,\quad d=1.
\]

Hledáme den, kdy průměr je 130 cm:
\[
a_n=130.
\]

Použijeme vzorec:
\[
a_n=a_1+(n-1)d.
\]

Dosadíme:
\[
130=51+(n-1).
\]

\[
130=50+n,
\]

\[
n=80.
\]

\bigskip

\textbf{Odpověď:}  
Průměr \(130\) cm mělo 80 dlaždic.

\section*{Úloha č. 9.2}

\textbf{Řešení:}

Každý den pokládají o jednu dlaždici více:
\[\
1,2,3,\dots,80.
\]

Celkový počet dlaždic je součet:
\[
S_n=\frac{n}{2}(a_1+a_n).
\]

Dosadíme:
\[
S_{80}=\frac{80}{2}(1+80)=40\cdot 81=3240.
\]

\bigskip

\textbf{Odpověď:}  
Celkem bylo položeno 3240 dlaždic.

\section*{Úloha č. 9.3}

\textbf{Řešení:}

Hledáme průměr dlaždice, která byla položena jako tisící v pořadí.

Nejprve zjistíme, ve kterém dni byla položena.

Hledáme nejmenší \(n\), pro které platí:
\[
\frac{n(n+1)}{2} \ge 1000.
\]

Zkusíme:
\[
44\cdot 45/2 = 990,
\]
\[
45\cdot 46/2 = 1035.
\]

Tisící dlaždice tedy byla položena v 45. dni.

\medskip

Průměr dlaždice v \(n\)-tý den je:
\[
a_n=51+(n-1).
\]

Pro \(n=45\):
\[
a_{45}=51+44=95.
\]

\bigskip

\textbf{Odpověď:}   
Tisící dlaždice měla průměr 95 cm.


% ---------------------------------
% ------ Ú L O H A  -  č. 10 ------
% ---------------------------------
\newpage
\section*{Úloha č. 10}

\textbf{Zadání:}  
Členové kapely si pořídili aparaturu za \(13\,500\) Kč. Všichni na nákup přispěli stejnou částkou. Kdyby jim vypomohl ještě fanoušek Jarda a celkovou částku si rovnocenně rozdělili i s ním, každému z členů kapely by se náklady snížily o \(450\) Kč.

Užitím rovnice nebo soustavy rovnic vypočtěte, kolik členů má kapela.

\bigskip

\textbf{Řešení:}

Označme si:
\begin{itemize}
    \item \(x\) \dots počet členů kapely.
\end{itemize}

Každý člen kapely původně zaplatil
\[
\frac{13500}{x}.
\]

Kdyby přispěl ještě fanoušek Jarda, dělila by se částka mezi
\[
x+1
\]
lidí. Každý by pak zaplatil
\[
\frac{13500}{x+1}.
\]

Podle zadání by se každému členovi kapely náklady snížily o \(450\) Kč. Platí tedy:
\[
\frac{13500}{x}-\frac{13500}{x+1}=450.
\]

\medskip

Rovnici upravíme na společný jmenovatel:
\[
\frac{13500(x+1)-13500x}{x(x+1)}=450.
\]

V čitateli se členy s \(x\) odečtou:
\[
\frac{13500}{x(x+1)}=450.
\]

Vynásobíme:
\[
13500=450x(x+1).
\]

Vydělíme číslem \(450\):
\[
30=x(x+1).
\]

Dostáváme rovnici:
\[
x^2+x-30=0.
\]

Rozložíme na součin:
\[
(x+6)(x-5)=0.
\]

Odtud:
\[
x=-6 \quad \text{nebo} \quad x=5.
\]

Záporné řešení nedává smysl, proto:
\[
x=5.
\]

\bigskip

\textbf{Kontrola:}

Pokud má kapela 5 členů, každý původně zaplatí
\[
\frac{13500}{5}=2700\text{ Kč}.
\]

Kdyby platil i Jarda, dělilo by se mezi 6 lidí:
\[
\frac{13500}{6}=2250\text{ Kč}.
\]

Rozdíl je
\[
2700-2250=450\text{ Kč},
\]
což souhlasí.

\bigskip

\textbf{Odpověď:}  
Kapela má 5 členů.


% ---------------------------------
% ------ Ú L O H A  -  č. 11 ------
% ---------------------------------
\newpage
\section*{Úloha č. 11}

\textbf{Zadání:}  
Hvězdičky jsou v obrazci umístěny v řadách nad sebou. Počty hvězdiček v jednotlivých řadách tvoří konečnou aritmetickou posloupnost. V nejkratší řadě jsou dvě hvězdičky. Počet hvězdiček v nejdelší řadě je o \(99\) větší než počet všech řad.

Kolik hvězdiček obsahuje celý obrazec?

\bigskip

\textbf{Řešení:}

Označme si:
\begin{itemize}
    \item \(n\) \dots počet řad,
    \item \(a_1=2\) \dots počet hvězdiček v nejkratší řadě,
    \item \(a_n\) \dots počet hvězdiček v nejdelší řadě.
\end{itemize}

Z obrázku je vidět, že počty hvězdiček v řadách tvoří aritmetickou posloupnost a mezi sousedními řadami přibývají vždy 3 hvězdičky. Tedy:
\[
d=3.
\]

Počet hvězdiček v poslední řadě je podle zadání o \(99\) větší než počet všech řad:
\[
a_n=n+99.
\]

Zároveň pro \(n\)-tý člen aritmetické posloupnosti platí:
\[
a_n=a_1+(n-1)d.
\]

Dosadíme:
\[
n+99=2+(n-1)\cdot 3.
\]

Upravíme pravou stranu:
\[
n+99=2+3n-3,
\]
\[
n+99=3n-1.
\]

Převedeme členy:
\[
100=2n,
\]
\[
n=50.
\]

Počet řad je tedy
\[
50.
\]

Nyní určíme počet hvězdiček v poslední řadě:
\[
a_n=n+99=50+99=149.
\]

Celkový počet hvězdiček je součet aritmetické posloupnosti:
\[
S_n=\frac{n}{2}(a_1+a_n).
\]

Dosadíme:
\[
S_{50}=\frac{50}{2}(2+149)=25\cdot 151=3775.
\]

\bigskip

\textbf{Odpověď:}  
Celý obrazec obsahuje 3775 hvězdiček.


% ---------------------------------
% ------ Ú L O H A  -  č. 12 ------
% ---------------------------------
\newpage
\section*{Úloha č. 12}

\textbf{Zadání:}  
Kladná čísla \(g, h\) se liší o 33. Číslo \(g\) je o \(20\%\) větší než neznámé číslo \(x\), neznámé číslo \(x\) je o \(20\%\) větší než číslo \(h\).

Užitím rovnice nebo soustavy rovnic vypočtěte neznámé číslo \(x\).

\bigskip

\textbf{Řešení:}

Ze zadání víme:
\[
g = 1{,}2x
\]
a zároveň
\[
x = 1{,}2h.
\]

Dosadíme druhý vztah do prvního:
\[
g = 1{,}2 \cdot 1{,}2h = 1{,}44h.
\]

Protože čísla \(g\) a \(h\) se liší o 33, platí:
\[
g-h=33.
\]

Dosadíme za \(g\):
\[
1{,}44h-h=33.
\]

\[
0{,}44h=33.
\]

\[
h=\frac{33}{0{,}44}=75.
\]

Nyní dopočítáme číslo \(x\):
\[
x=1{,}2h=1{,}2\cdot 75=90.
\]

\bigskip

\textbf{Kontrola:}

\[
g=1{,}2x=1{,}2\cdot 90=108.
\]

Rozdíl:
\[
108-75=33,
\]
což souhlasí.

\bigskip

\textbf{Odpověď:}  
Neznámé číslo je 90.


% ---------------------------------
% ------ Ú L O H A  -  č. 13 ------
% ---------------------------------
\newpage
\section*{Úloha č. 13}

\textbf{Zadání:}  
Do rovnoramenného trojúhelníku \(ABC\) je vepsáno nekonečně mnoho čtverců. Jedna strana každého čtverce leží na základně \(AB\) trojúhelníku. Čtverce se vzájemně dotýkají.

Největší čtverec s délkou strany \(20\ \mathrm{mm}\) je umístěn tak, že osa trojúhelníku je současně osou čtverce. Každé dva sousední čtverce mají jeden společný vrchol a délky jejich stran jsou v poměru \(5:4\).

Vypočtěte v \(\mathrm{mm}^2\) obsah trojúhelníku \(ABC\).

\bigskip

\textbf{Řešení:}

Největší čtverec má stranu
\[
20\ \mathrm{mm}.
\]

Sousledné čtverce mají délky stran v poměru
\[
5:4.
\]
Proto čtverec sousedící s největším čtvercem má stranu
\[
20\cdot \frac{4}{5}=16\ \mathrm{mm}.
\]

\medskip

Budeme uvažovat pravou polovinu trojúhelníku.

Označme:
\begin{itemize}
    \item \(D\) \dots pravý horní vrchol největšího čtverce,
    \item \(E\) \dots pravý horní vrchol sousedního menšího čtverce.
\end{itemize}

Protože oba tyto vrcholy leží na rameni trojúhelníku, body \(D\) a \(E\) leží na téže přímce.

\medskip

Vodorovná vzdálenost bodů \(D\) a \(E\) je rovna straně menšího čtverce:
\[
16\ \mathrm{mm}.
\]

Svislá vzdálenost těchto bodů je rozdíl délek stran obou čtverců:
\[
20-16=4\ \mathrm{mm}.
\]

Rameno trojúhelníku tedy klesá o \(4\ \mathrm{mm}\) na každých \(16\ \mathrm{mm}\) vodorovného posunu. Poměr svislé a vodorovné změny je tedy
\[
4:16=1:4.
\]

\bigskip

\textbf{Výška trojúhelníku}

Největší čtverec je souměrně umístěn podle osy trojúhelníku, takže jeho pravý horní vrchol je od osy vzdálen
\[
10\ \mathrm{mm}.
\]

Jestliže na každé \(4\ \mathrm{mm}\) vodorovně připadá změna výšky o \(1\ \mathrm{mm}\), pak při vodorovné vzdálenosti \(10\ \mathrm{mm}\) odpovídá svislá změna
\[
\frac{10}{4}=2{,}5\ \mathrm{mm}.
\]

Výška trojúhelníku je tedy
\[
v = 20+2{,}5=22{,}5\ \mathrm{mm}.
\]

\bigskip

\textbf{Délka základny}

Od pravého horního vrcholu největšího čtverce do pravého krajního bodu základny klesne rameno trojúhelníku z výšky \(20\ \mathrm{mm}\) na nulu.

Při poměru \(1:4\) tomu odpovídá vodorovná vzdálenost
\[
20\cdot 4=80\ \mathrm{mm}.
\]

Od osy trojúhelníku k pravému hornímu vrcholu čtverce je ještě
\[
10\ \mathrm{mm},
\]
takže polovina základny trojúhelníku je
\[
10+80=90\ \mathrm{mm}.
\]

Celá základna tedy měří
\[
AB=180\ \mathrm{mm}.
\]

\bigskip

\textbf{Obsah trojúhelníku}

Obsah trojúhelníku je
\[
S=\frac{AB\cdot v}{2}.
\]

Dosadíme:
\[
S=\frac{180\cdot 22{,}5}{2}=90\cdot 22{,}5=2025\ \mathrm{mm}^2.
\]

\bigskip

\textbf{Odpověď:}  
Obsah trojúhelníku \(ABC\) je 2025 \(\mathrm{mm}^2\).


% ---------------------------------
% ------ Ú L O H A  -  č. 14 ------
% ---------------------------------
\newpage
\section*{Úloha č. 14}

\textbf{Zadání:}  
Adam, Bořek a Cyril si koupili ze společných příspěvků doplňky k hudební aparatuře.

Cyril s Bořkem přispěli dohromady částkou \(5\,100\) korun. Bořek přispěl o třetinu vyšší částkou než Adam a Adam částkou o třetinu nižší než Cyril.

Po zakoupení všech doplňků chlapci věnovali zbývající částku na dobročinné účely. Její hodnota se rovnala pětině částky, kterou zaplatili za všechny doplňky.

Užitím rovnice nebo soustavy rovnic vypočtěte v korunách
\begin{enumerate}
    \item[4.1] příspěvek Adama,
    \item[4.2] částku věnovanou na dobročinné účely.
\end{enumerate}

\bigskip

\textbf{Řešení:}

Označme si:
\begin{itemize}
    \item \(A\) \dots Adamův příspěvek,
    \item \(B\) \dots Bořkův příspěvek,
    \item \(C\) \dots Cyrilův příspěvek.
\end{itemize}

Ze zadání víme:
\[
B+C=5100.
\]

Dále Bořek přispěl o třetinu vyšší částkou než Adam, tedy
\[
B=\frac{4}{3}A.
\]

Adam přispěl o třetinu nižší částkou než Cyril, tedy
\[
A=\frac{2}{3}C.
\]

Z poslední rovnice vyjádříme:
\[
C=\frac{3}{2}A.
\]

Dosadíme do rovnice \(B+C=5100\):
\[
\frac{4}{3}A+\frac{3}{2}A=5100.
\]

Převedeme na společného jmenovatele:
\[
\frac{8}{6}A+\frac{9}{6}A=5100,
\]
\[
\frac{17}{6}A=5100.
\]

Odtud:
\[
A=5100\cdot \frac{6}{17}=1800.
\]

Adamův příspěvek je tedy
\[
1800\text{ Kč}.
\]

\medskip

Nyní dopočítáme ostatní příspěvky:
\[
B=\frac{4}{3}\cdot 1800=2400,
\]
\[
C=\frac{3}{2}\cdot 1800=2700.
\]

Celkem chlapci přispěli:
\[
1800+2400+2700=6900\text{ Kč}.
\]

Označme si:
\begin{itemize}
    \item \(x\) \dots cena všech doplňků,
    \item \(y\) \dots částka věnovaná na dobročinné účely.
\end{itemize}

Podle zadání platí:
\[
y=\frac{1}{5}x
\]
a zároveň
\[
x+y=6900.
\]

Dosadíme:
\[
x+\frac{1}{5}x=6900,
\]
\[
\frac{6}{5}x=6900,
\]
\[
x=6900\cdot \frac{5}{6}=5750.
\]

Potom
\[
y=6900-5750=1150.
\]

\bigskip

\textbf{Odpověď:}
\begin{itemize}
    \item[14.1] Adam přispěl 1800 Kč.
    \item[14.2] Na dobročinné účely věnovali 1150 Kč.
\end{itemize}



% ---------------------------------
% ------ Ú L O H A  -  č. 15 ------
% ---------------------------------
\newpage
\section*{Úloha č. 15}

\textbf{Zadání:}  
Každý pracovník firmy mluví alespoň jedním ze dvou jazyků – anglicky nebo německy.

Polovina těch, kteří mluví německy, mluví i anglicky.  
Třetina těch, kteří mluví anglicky, mluví i německy.

Vyjádřete zlomkem v základním tvaru, jaká část všech pracovníků firmy mluví oběma jazyky (anglicky i německy).

\bigskip

\textbf{Řešení:}

Označme si:
\begin{itemize}
    \item \(N\) \dots počet pracovníků, kteří mluví německy,
    \item \(A\) \dots počet pracovníků, kteří mluví anglicky,
    \item \(x\) \dots počet pracovníků, kteří mluví oběma jazyky.
\end{itemize}

Ze zadání víme:
\[
x=\frac{1}{2}N
\]
a zároveň
\[
x=\frac{1}{3}A.
\]

Odtud:
\[
N=2x,\qquad A=3x.
\]

Všichni pracovníci mluví alespoň jedním jazykem, takže celkový počet pracovníků je
\[
N+A-x.
\]

Dosadíme:
\[
2x+3x-x=4x.
\]

Počet pracovníků, kteří mluví oběma jazyky, je \(x\), z celkového počtu \(4x\). Hledaný podíl je tedy
\[
\frac{x}{4x}=\frac{1}{4}.
\]

\bigskip

\textbf{Odpověď:}  
Oběma jazyky mluví \(\frac{1}{4}\) všech pracovníků firmy.



% ---------------------------------
% ------ Ú L O H A  -  č. 16 ------
% ---------------------------------
\newpage
\section*{Úloha č. 16}

\textbf{Zadání:}  
Třetí mocnina neznámého čísla je o 100 menší než druhá mocnina součtu téhož neznámého čísla s číslem 10.

Vypočtěte neznámé číslo. Najděte všechna řešení.

\bigskip

\textbf{Řešení:}

Označme neznámé číslo \(x\).

Podle zadání platí:
\[
x^3=(x+10)^2-100.
\]

Nejprve upravíme pravou stranu:
\[
(x+10)^2-100=x^2+20x+100-100=x^2+20x.
\]

Dostáváme rovnici:
\[
x^3=x^2+20x.
\]

Vše převedeme na jednu stranu:
\[
x^3-x^2-20x=0.
\]

Vytkneme \(x\):
\[
x(x^2-x-20)=0.
\]

Kvadratický trojčlen rozložíme:
\[
x(x-5)(x+4)=0.
\]

Odtud dostáváme tři řešení:
\[
x=0,\qquad x=5,\qquad x=-4.
\]

\bigskip

\textbf{Odpověď:}  
Neznámé číslo může být -4,\ 0,\ 5.


% ---------------------------------
% ------ Ú L O H A  -  č. 17 ------
% ---------------------------------
\newpage
\section*{Úloha č. 17}

\textbf{Zadání:}  
K dispozici je 11 ampulí o stejném objemu.

Do 3 ampulí dohromady se vejde více než 4 centilitry tekutiny, ale méně než 5 centilitrů.  
Do 4 ampulí dohromady se vejde více než 5 centilitrů tekutiny, ale méně než 6 centilitrů.  
Celkový objem 11 ampulí je přesně \(n\) centilitrů, kde \(n\) je přirozené číslo.

Vypočtěte \(n\).

\bigskip

\textbf{Řešení:}

Označme objem jedné ampule \(v\) centilitrů.

Ze zadání plyne:
\[
4<3v<5
\]
a zároveň
\[
5<4v<6.
\]

Z první nerovnosti dostaneme po vydělení třemi:
\[
\frac{4}{3}<v<\frac{5}{3}.
\]

Z druhé nerovnosti dostaneme po vydělení čtyřmi:
\[
\frac{5}{4}<v<\frac{6}{4}=\frac{3}{2}.
\]

Objem jedné ampule tedy musí splňovat obě podmínky současně:
\[
\frac{4}{3}<v<\frac{3}{2}.
\]

Hledáme takové \(v\), aby celkový objem 11 ampulí byl přirozené číslo:
\[
11v=n.
\]

Protože
\[
\frac{4}{3}<v<\frac{3}{2},
\]
po vynásobení číslem 11 dostaneme:
\[
\frac{44}{3}<11v<\frac{33}{2}.
\]

Tedy:
\[
14{,}6\overline{6}<n<16{,}5.
\]

Jediné přirozené číslo v tomto intervalu je
\[
n=15.
\]

\bigskip

\textbf{Odpověď:}  
\[
n=15
\].


% ---------------------------------
% ------ Ú L O H A  -  č. 18 ------
% ---------------------------------
\newpage
\section*{Úloha č. 18}

\textbf{Zadání:}  
Byla vybrána skupina tří zaměstnanců, jejichž platy tvoří tři po sobě jdoucí členy aritmetické posloupnosti. Z těchto tří platů je nejvyšší plat o \(25\ \%\) vyšší než nejnižší plat a aritmetický průměr všech tří platů je \(36\,000\) korun.

Vypočtěte, o kolik korun se liší nejvyšší a nejnižší plat ve vybrané skupině.

\bigskip

\textbf{Řešení:}

Tři po sobě jdoucí členy aritmetické posloupnosti si označíme:
\[
a-d,\quad a,\quad a+d.
\]

Protože aritmetický průměr těchto tří čísel je roven prostřednímu členu, platí:
\[
a=36\,000.
\]

Platy tedy jsou:
\[
36\,000-d,\quad 36\,000,\quad 36\,000+d.
\]

Ze zadání dále víme, že nejvyšší plat je o \(25\%\) vyšší než nejnižší plat. To znamená:
\[
36\,000+d = 1{,}25(36\,000-d).
\]

Rovnici upravíme:
\[
36\,000+d = 45\,000-1{,}25d.
\]

\[
d+1{,}25d = 45\,000-36\,000.
\]

\[
2{,}25d = 9\,000.
\]

\[
d = 4\,000.
\]

Nejnižší a nejvyšší plat se tedy liší o
\[
(36\,000+d)-(36\,000-d)=2d=8\,000.
\]

\bigskip

\textbf{Odpověď:}  
Nejvyšší a nejnižší plat se liší o \(8\,000\text{ Kč}\).


% ---------------------------------
% ------ Ú L O H A  -  č. 19 ------
% ---------------------------------
\newpage
\section*{Úloha č. 19.1}

\textbf{Zadání:}  
Každý ze zobrazených čtverců má dva sousední vrcholy na přímce \(AE\), další vrchol na přímce \(CE\) a poslední vrchol na přímce \(DE\).

Zobrazené čtverce splňují současně následující podmínky:
\begin{itemize}
    \item délka strany prvního čtverce je \(a_1=1\ \mathrm{dm}\),
    \item délky stran všech po sobě jdoucích čtverců tvoří geometrickou posloupnost,
    \item délka strany každého čtverce je součtem délek stran dvou následujících čtverců.
\end{itemize}

Vyjádřete (bez zaokrouhlení) kvocient \(q\) dané posloupnosti.

\bigskip

\textbf{Řešení:}

Označme délky stran čtverců:
\[
a_1,\ a_2,\ a_3,\dots
\]

Protože tvoří geometrickou posloupnost, platí:
\[
a_2=a_1q,\qquad a_3=a_1q^2.
\]

Podle zadání je délka strany každého čtverce součtem délek stran dvou následujících čtverců, tedy například:
\[
a_1=a_2+a_3.
\]

Dosadíme:
\[
a_1=a_1q+a_1q^2.
\]

Vytkneme \(a_1\):
\[
a_1(1-q-q^2)=0.
\]

Protože \(a_1\neq 0\), musí platit:
\[
1-q-q^2=0.
\]

Upravíme:
\[
q^2+q-1=0.
\]

Použijeme kvadratický vzorec:
\[
q=\frac{-1\pm\sqrt{1+4}}{2}=\frac{-1\pm\sqrt{5}}{2}.
\]

Protože délky stran jsou kladné, musí být i kvocient kladný. Proto:
\[
q=\frac{-1+\sqrt{5}}{2}.
\]

\bigskip

\textbf{Odpověď:}  
\[
q=\frac{\sqrt{5}-1}{2}.
\]


\section*{Úloha č. 19.2}

\textbf{Zadání:}  
Vypočtěte v dm (bez zaokrouhlení) délku úsečky \(AE\).

\bigskip

\textbf{Řešení:}

Délka úsečky \(AE\) je součtem délek stran všech čtverců:
\[
AE=a_1+a_2+a_3+\dots
\]

Jde o nekonečnou geometrickou řadu s prvním členem
\[
a_1=1
\]
a kvocientem
\[
q=\frac{\sqrt{5}-1}{2}.
\]

Protože \(|q|<1\), můžeme použít vzorec pro součet nekonečné geometrické řady:
\[
S=\frac{a_1}{1-q}.
\]

Dosadíme:
\[
AE=\frac{1}{1-\frac{\sqrt{5}-1}{2}}.
\]

Upravíme jmenovatel:
\[
1-\frac{\sqrt{5}-1}{2}=\frac{2-(\sqrt{5}-1)}{2}=\frac{3-\sqrt{5}}{2}.
\]

Tedy:
\[
AE=\frac{1}{\frac{3-\sqrt{5}}{2}}=\frac{2}{3-\sqrt{5}}.
\]

Odstraníme odmocninu ze jmenovatele:
\[
AE=\frac{2}{3-\sqrt{5}}\cdot\frac{3+\sqrt{5}}{3+\sqrt{5}}
=\frac{2(3+\sqrt{5})}{9-5}
=\frac{2(3+\sqrt{5})}{4}
=\frac{3+\sqrt{5}}{2}.
\]

\bigskip

\textbf{Odpověď:}  
\[
AE=\frac{3+\sqrt{5}}{2}\ \mathrm{dm}.
\]


\section*{Úloha č. 19.3}

\textbf{Zadání:}  
Vypočtěte v \(\mathrm{dm}^2\) obsah trojúhelníku \(CDE\).

\bigskip

\textbf{Řešení:}

Z obrázku je vidět, že:
\begin{itemize}
    \item bod \(D\) leží svisle nad bodem \(A\),
    \item bod \(C\) leží svisle nad bodem \(B\),
    \item strana prvního čtverce má délku \(a_1=1\ \mathrm{dm}\).
\end{itemize}

Proto:
\[
DC=AB=a_1=1\ \mathrm{dm}.
\]

Výška trojúhelníku \(CDE\) na základnu \(DC\) je vzdálenost bodu \(E\) od přímky \(DC\).  
Přímka \(DC\) je vodorovná ve výšce \(1\ \mathrm{dm}\) nad přímkou \(AE\), na které leží bod \(E\). Výška trojúhelníku je tedy:
\[
v=1\ \mathrm{dm}.
\]

Obsah trojúhelníku:
\[
S=\frac{DC\cdot v}{2}=\frac{1\cdot 1}{2}=\frac12.
\]

\bigskip

\textbf{Odpověď:}  
\[
\frac12\ \mathrm{dm}^2.
\]



% ---------------------------------
% ------ Ú L O H A  -  č. 20 ------
% ---------------------------------
\newpage
\section*{Úloha č. 20.1}

\textbf{Zadání:}  
V dílně se pracuje na dvou strojích.

Na prvním stroji se za určitou dobu vyrobí \(n\) výrobků, zatímco na druhém stroji se za dobu o třetinu delší vyrobí \(2n\) výrobků.

Vypočtěte, kolikrát více výrobků se za stejnou dobu vyrobí na druhém stroji než na prvním stroji.

\bigskip

\textbf{Řešení:}

Označme dobu, za kterou první stroj vyrobí \(n\) výrobků, jako \(t\).

Pak výkon prvního stroje je:
\[
\frac{n}{t}.
\]

Druhý stroj vyrobí \(2n\) výrobků za dobu o třetinu delší, tedy za čas
\[
\frac{4}{3}t.
\]

Jeho výkon je:
\[
\frac{2n}{\frac{4}{3}t}=2n\cdot\frac{3}{4t}=\frac{3n}{2t}.
\]

Hledáme, kolikrát je výkon druhého stroje větší než výkon prvního:
\[
\frac{\frac{3n}{2t}}{\frac{n}{t}}=\frac{3}{2}.
\]

\bigskip

\textbf{Odpověď:}  
Na druhém stroji se za stejnou dobu vyrobí \(\frac{3}{2}\)-krát více výrobků než na prvním stroji.


\section*{Úloha č. 20.2}

\textbf{Zadání:}  
Na obou strojích se za \(p\) hodin vyrobí dohromady \(10n\) výrobků.

V závislosti na \(p\) vyjádřete, za kolik hodin se \(10n\) výrobků vyrobí jen na prvním stroji.

\bigskip

\textbf{Řešení:}

Z předchozí části víme, že výkon druhého stroje je
\[
\frac32
\]
násobkem výkonu prvního stroje.

Označme výkon prvního stroje jako
\[
v.
\]
Pak výkon druhého stroje je
\[
\frac32v.
\]

Společný výkon obou strojů je:
\[
v+\frac32v=\frac52v.
\]

Podle zadání se za \(p\) hodin vyrobí dohromady \(10n\) výrobků, tedy společný výkon odpovídá hodnotě
\[
\frac{10n}{p}.
\]

Platí tedy:
\[
\frac52v=\frac{10n}{p}.
\]

Odtud:
\[
v=\frac{10n}{p}\cdot\frac{2}{5}=\frac{4n}{p}.
\]

První stroj tedy vyrábí rychlostí
\[
\frac{4n}{p}
\]
výrobků za hodinu.

Hledáme čas \(T\), za který první stroj sám vyrobí \(10n\) výrobků:
\[
\frac{4n}{p}\cdot T=10n.
\]

Odtud:
\[
T=10n\cdot\frac{p}{4n}=\frac{10p}{4}=\frac{5p}{2}.
\]

\bigskip

\textbf{Odpověď:}  
Jen na prvním stroji se \(10n\) výrobků vyrobí za \(\frac{5p}{2}\) hodin.



% ---------------------------------
% ------ Ú L O H A  -  č. 21 ------
% ---------------------------------
\newpage
\section*{Úloha č. 21}

\textbf{Zadání:}  
Každý z 10 soutěžících mohl získat 0 až 3 trestné body.

Četnosti představující počty soutěžících, kteří získali stejný počet trestných bodů, jsou v první tabulce uvedeny v nesprávném pořadí.

Při správném přiřazení četností je medián a modus počtu trestných bodů stejný.  
Ze všech takových možností je správná pouze ta, která vede k nejmenší možné hodnotě aritmetického průměru počtu trestných bodů.

Z údajů ve správně vyplněné tabulce vypočtěte aritmetický průměr počtu trestných bodů.

\bigskip

\textbf{Řešení:}

Četnosti jsou čísla
\[
0,\ 2,\ 3,\ 5,
\]
ale nejsou správně přiřazena k počtům trestných bodů
\[
0,\ 1,\ 2,\ 3.
\]

Protože největší četnost je \(5\), bude modus právě ten počet bodů, ke kterému přiřadíme četnost \(5\).

\medskip

Aby byl medián stejný jako modus, nemůže být modus:
\begin{itemize}
    \item \(0\), protože pak by prvních 5 hodnot bylo \(0\), ale šestá už by byla větší než 0,
    \item \(3\), protože pak by posledních 5 hodnot bylo \(3\), ale pátá by byla menší než 3.
\end{itemize}

Modus tedy může být jen
\[
1 \quad \text{nebo} \quad 2.
\]

\medskip

Ze všech takových možností máme vybrat tu, která dává \textbf{nejmenší aritmetický průměr}.  
Aby byl průměr co nejmenší, je výhodné přiřadit co největší četnosti co nejmenším počtům bodů.

Proto zvolíme:
\[
\begin{array}{c|cccc}
\text{Trestné body} & 0 & 1 & 2 & 3 \\
\hline
\text{Četnost} & 3 & 5 & 2 & 0
\end{array}
\]

Zkontrolujeme:
\begin{itemize}
    \item modus je \(1\), protože má největší četnost \(5\),
    \item uspořádaný soubor je
    \[
    0,0,0,1,1,1,1,1,2,2,
    \]
    takže medián je průměr 5. a 6. hodnoty:
    \[
    \frac{1+1}{2}=1.
    \]
\end{itemize}

Medián i modus jsou tedy opravdu stejné.

\medskip

Nyní spočítáme aritmetický průměr:
\[
\overline{x}=\frac{0\cdot 3+1\cdot 5+2\cdot 2+3\cdot 0}{10}
=\frac{0+5+4+0}{10}
=\frac{9}{10}=0{,}9.
\]

\bigskip

\textbf{Odpověď:}  
Aritmetický průměr počtu trestných bodů je \(0{,}9\).



% ---------------------------------
% ------ Ú L O H A  -  č. 22 ------
% ---------------------------------
\newpage
\section*{Úloha č. 22.1}

\textbf{Zadání:}  
Šestiúhelník \(ABCDEF\) je rozdělen na dva bílé rovnoramenné trojúhelníky a dva shodné tmavé čtverce. Všechny čtyři útvary mají společný vrchol \(V\).

Součet obsahů obou bílých trojúhelníků označme \(S_B\).  
Součet obsahů obou tmavých čtverců označme \(S_T\).  
Úhel \(DVE\) má velikost \(\varphi \in \left(0;\frac{\pi}{2}\right)\).

Vyjádřete poměr obsahů \(S_B : S_T\) v závislosti na velikosti úhlu \(\varphi\).

\bigskip

\textbf{Řešení:}

Z obrázku je patrné, že:
\begin{itemize}
    \item oba bílé útvary jsou shodné rovnoramenné trojúhelníky,
    \item oba tmavé útvary jsou shodné čtverce,
    \item délky stran označené \(a\) jsou stranami těchto útvarů.
\end{itemize}

\medskip

\textbf{Obsah jednoho bílého trojúhelníku}

Každý bílý trojúhelník má dvě strany délky \(a\) a úhel mezi nimi \(\varphi\).  
Jeho obsah tedy je:
\[
S_{\triangle}=\frac12 a^2 \sin\varphi.
\]

Jsou dva takové trojúhelníky, proto:
\[
S_B=2\cdot \frac12 a^2\sin\varphi=a^2\sin\varphi.
\]

\medskip

\textbf{Obsah tmavých čtverců}

Jeden čtverec se stranou \(a\) má obsah
\[
a^2.
\]

Jsou dva, tedy:
\[
S_T=2a^2.
\]

\medskip

Poměr obsahů je:
\[
S_B:S_T = a^2\sin\varphi : 2a^2.
\]

Po krácení \(a^2\) dostáváme:
\[
S_B:S_T = \sin\varphi : 2.
\]

\bigskip

\textbf{Odpověď:}  
\[
S_B:S_T=\sin\varphi:2.
\]


\section*{Úloha č. 22.2}

\textbf{Zadání:}  
Vypočtěte velikost úhlu \(\varphi\), je-li poměr obsahů \(S_B:S_T\) roven \(1:4\).

\bigskip

\textbf{Řešení:}

Z předchozí části víme, že
\[
S_B:S_T=\sin\varphi:2.
\]

Podle zadání platí:
\[
\sin\varphi:2 = 1:4.
\]

To znamená:
\[
\frac{\sin\varphi}{2}=\frac14.
\]

Odtud:
\[
\sin\varphi=\frac12.
\]

Protože
\[
\varphi \in \left(0;\frac{\pi}{2}\right),
\]
připadá v úvahu jen řešení
\[
\varphi=\frac{\pi}{6}.
\]

\bigskip

\textbf{Odpověď:}  
\[
\boxed{\varphi=\frac{\pi}{6}}
\]



% ---------------------------------
% ------ Ú L O H A  -  č. 23 ------
% ---------------------------------
\newpage
\section*{Úloha č. 23}

\textbf{Zadání:}  
Ve čtverci \(PQRS\) o straně délky \(6\ \mathrm{m}\) je nekonečně mnoho stále se zmenšujících tmavých rovnoramenných pravoúhlých trojúhelníků. Největší z nich je trojúhelník \(PQR\).

Každý následující trojúhelník má vrchol pravého úhlu uprostřed přepony předchozího trojúhelníku, což je i jediný společný bod dvou trojúhelníků.

Středem stejnolehlosti libovolné dvojice těchto trojúhelníků je vrchol \(S\).

Vyjádřete poměr obsahu všech bílých ploch ku obsahu všech tmavých ploch ve čtverci \(PQRS\).

\bigskip

\textbf{Řešení:}

Obsah celého čtverce je
\[
S_{\square}=6^2=36\ \mathrm{m}^2.
\]

\medskip

\textbf{Obsah největšího tmavého trojúhelníku}

Trojúhelník \(PQR\) je pravoúhlý rovnoramenný s odvěsnami délky \(6\ \mathrm{m}\), takže jeho obsah je
\[
S_1=\frac12\cdot 6\cdot 6=18\ \mathrm{m}^2.
\]

\medskip

\textbf{Poměr obsahů po sobě jdoucích trojúhelníků}

Zadání říká, že všechny trojúhelníky jsou si podobné a středem stejnolehlosti je bod \(S\).  
Každý další trojúhelník má vrchol pravého úhlu ve středu přepony předchozího trojúhelníku, proto se všechny délky zmenší na polovinu.

Lineární koeficient podobnosti je tedy
\[
\frac12.
\]

Obsahy se proto zmenšují v poměru
\[
\left(\frac12\right)^2=\frac14.
\]

Obsahy tmavých trojúhelníků tedy tvoří geometrickou řadu:
\[
18,\ 18\cdot\frac14,\ 18\cdot\frac{1}{4^2},\dots
\]

Součet obsahů všech tmavých trojúhelníků je:
\[
S_T=\frac{18}{1-\frac14}
=\frac{18}{\frac34}
=24\ \mathrm{m}^2.
\]

\medskip

Obsah všech bílých ploch je tedy:
\[
S_B=36-24=12\ \mathrm{m}^2.
\]

Hledaný poměr:
\[
S_B:S_T=12:24=1:2.
\]

\bigskip

\textbf{Odpověď:}  
Poměr obsahu všech bílých ploch ku obsahu všech tmavých ploch je
\[
1:2.
\]



% ---------------------------------
% ------ Ú L O H A  -  č. 24 ------
% ---------------------------------
\newpage
\section*{Úloha č. 24.1}

\textbf{Zadání:}  
Z dlouhodobých statistik vyplývá, že \(5\%\) osob má vlastnost \(A\) a \(10\%\) osob vlastnost \(B\). Obě vlastnosti se nepodmiňují, naopak jsou nezávislé.

Vypočtěte s přesností na tisíciny pravděpodobnost jevu:  
Náhodně vybraná osoba nemá ani jednu z obou vlastností \(A, B\).

\bigskip

\textbf{Řešení:}

Označme si:
\[
P(A)=0{,}05,\qquad P(B)=0{,}10.
\]

Hledáme pravděpodobnost, že osoba nemá ani jednu z vlastností:
\[
P(\overline{A}\cap \overline{B}).
\]

Protože jsou jevy \(A\) a \(B\) nezávislé, jsou nezávislé i jejich doplňky. Platí tedy:
\[
P(\overline{A}\cap \overline{B})=P(\overline{A})\cdot P(\overline{B}).
\]

Spočítáme doplňkové pravděpodobnosti:
\[
P(\overline{A})=1-0{,}05=0{,}95,
\]
\[
P(\overline{B})=1-0{,}10=0{,}90.
\]

Dosadíme:
\[
P(\overline{A}\cap \overline{B})=0{,}95\cdot 0{,}90=0{,}855.
\]

\bigskip

\textbf{Odpověď:}  
Hledaná pravděpodobnost je \(0{,}855\).


\section*{Úloha č. 24.2}

\textbf{Zadání:}  
Z dlouhodobých statistik vyplývá, že \(5\%\) osob má vlastnost \(A\) a \(10\%\) osob vlastnost \(B\). Obě vlastnosti se nepodmiňují, naopak jsou nezávislé.

Vypočtěte s přesností na tisíciny pravděpodobnost jevu:  
Alespoň jedna z 5 náhodně vybraných osob má alespoň jednu z obou vlastností \(A, B\).

\bigskip

\textbf{Řešení:}

Nejprve spočítáme pravděpodobnost, že jedna osoba má alespoň jednu z vlastností \(A\) nebo \(B\).

Platí:
\[
P(A\cup B)=P(A)+P(B)-P(A\cap B).
\]

Protože jsou jevy \(A\) a \(B\) nezávislé:
\[
P(A\cap B)=P(A)\cdot P(B)=0{,}05\cdot 0{,}10=0{,}005.
\]

Dosadíme:
\[
P(A\cup B)=0{,}05+0{,}10-0{,}005=0{,}145.
\]

Pravděpodobnost, že jedna osoba nemá ani jednu z těchto vlastností, je tedy:
\[
1-0{,}145=0{,}855.
\]

Hledáme pravděpodobnost, že \textbf{alespoň jedna z 5 osob} má alespoň jednu z vlastností.  
To je výhodnější spočítat přes doplněk:

\[
P(\text{alespoň jedna})=1-P(\text{žádná z 5 osob nemá ani jednu vlastnost}).
\]

Pravděpodobnost, že žádná z 5 osob nemá ani jednu vlastnost, je:
\[
0{,}855^5.
\]

Proto:
\[
P(\text{alespoň jedna})=1-0{,}855^5.
\]

Vypočítáme:
\[
0{,}855^5 \approx 0{,}457,
\]
tedy
\[
1-0{,}855^5 \approx 1-0{,}457=0{,}543.
\]

\bigskip

\textbf{Odpověď:}  
Hledaná pravděpodobnost je \(0{,}543\).




\section*{Úloha č. 24.3}

\textbf{Zadání:}  
Pětimístný kód má být sestaven z 5 různých číslic.

Z číslic \(0\)–\(9\) kód musí obsahovat číslice \(0,1,2\) a s nimi ještě další dvě číslice. Číslice \(0\) bude uprostřed kódu.

Jaký je počet všech možností pro sestavení kódu?

\bigskip

\textbf{Řešení:}

Kód má 5 míst a uprostřed musí být číslice
\[
0.
\]

Třetí pozice je tedy pevně daná.

Kód musí dále obsahovat číslice
\[
1 \quad \text{a} \quad 2
\]
a ještě dvě další různé číslice z množiny
\[
\{3,4,5,6,7,8,9\}.
\]

\medskip

\textbf{1. krok: výběr dvou dalších číslic}

Z číslic \(3\) až \(9\) vybíráme dvě:
\[
\binom{7}{2}=21.
\]

\medskip

\textbf{2. krok: rozmístění číslic}

Máme tedy k dispozici čtyři číslice:
\[
1,\ 2,\ a,\ b,
\]
které rozmísťujeme na zbývající čtyři pozice (první, druhou, čtvrtou a pátou).

Počet možností je:
\[
4! = 24.
\]

\medskip

Celkový počet kódů je:
\[
21\cdot 24=504.
\]

\bigskip

\textbf{Odpověď:}  
Počet všech možných kódů je \(504\).


% ---------------------------------
% ------ Ú L O H A  -  č. 25 ------
% ---------------------------------
\newpage
\section*{Úloha č. 25}

\textbf{Zadání:}  
Výrobní družstvo vyrobilo sadu stejných sekaček, každou s týmiž náklady na její výrobu.

Polovinu sekaček družstvo prodalo za cenu o \(70\%\) vyšší, než byly náklady na jejich výrobu.  
Každou další sekačku družstvo prodává za cenu o \(50\%\) vyšší, než byly náklady na její výrobu.

Přestože družstvo ještě neprodalo všechny sekačky, peníze získané za prodané sekačky již nyní přesně pokryly náklady na výrobu celé sady sekaček.

Vypočtěte, kolik procent všech vyrobených sekaček již družstvo prodalo.

\bigskip

\textbf{Řešení:}

Označme si:
\begin{itemize}
    \item \(n\) \dots celkový počet vyrobených sekaček,
    \item \(x\) \dots náklady na výrobu jedné sekačky.
\end{itemize}

Celkové náklady na výrobu celé sady jsou:
\[
nx.
\]

\medskip

Polovina sekaček, tedy
\[
\frac{n}{2},
\]
byla prodána za cenu o \(70\%\) vyšší než výrobní náklady, tedy za
\[
1{,}7x
\]
za kus.

Tržba z této poloviny je:
\[
\frac{n}{2}\cdot 1{,}7x = 0{,}85nx.
\]

\medskip

Označme si dále \(k\) počet dalších prodaných sekaček.  
Každá z nich byla prodána za cenu o \(50\%\) vyšší než výrobní náklady, tedy za
\[
1{,}5x.
\]

Tržba z těchto dalších prodaných sekaček je:
\[
1{,}5kx.
\]

\medskip

Podle zadání už celková tržba přesně pokryla náklady na výrobu celé sady, tedy:
\[
0{,}85nx + 1{,}5kx = nx.
\]

Odečteme \(0{,}85nx\):
\[
1{,}5kx = 0{,}15nx.
\]

Vydělíme \(x\):
\[
1{,}5k = 0{,}15n.
\]

Odtud:
\[
k = \frac{0{,}15}{1{,}5}n = 0{,}1n.
\]

Družstvo tedy kromě poloviny sekaček prodalo ještě dalších
\[
0{,}1n
\]
sekaček.

Celkem prodalo:
\[
\frac{n}{2}+0{,}1n = 0{,}5n+0{,}1n = 0{,}6n.
\]

To je
\[
60\%
\]
všech vyrobených sekaček.

\bigskip

\textbf{Odpověď:}  
Družstvo již prodalo \(60\%\) všech vyrobených sekaček.


% ---------------------------------
% ------ Ú L O H A  -  č. 26 ------
% ---------------------------------
\newpage
\section*{Úloha č. 26.1}

\textbf{Zadání:}  
Na obrázku jsou nakresleny tři obrazce složené ze čtverců o straně délky \(1\ \mathrm{cm}\).

Za těmito obrazci následují další. Obrazec s pořadovým číslem \(n>1\) se vytvoří z předchozího obrazce přidáním spodní řady, která obsahuje \(n\) čtverců.

Jeden z obrazců je složen z \(210\) čtverců.

Vypočtěte, kolikátý je tento obrazec.

\bigskip

\textbf{Řešení:}

Počty čtverců v jednotlivých obrazcích jsou:
\[
1,\ 1+2,\ 1+2+3,\dots
\]

\(n\)-tý obrazec tedy obsahuje součet prvních \(n\) přirozených čísel:
\[
S_n = 1+2+3+\dots+n = \frac{n(n+1)}{2}.
\]

Podle zadání platí:
\[
\frac{n(n+1)}{2}=210.
\]

Vynásobíme dvěma:
\[
n(n+1)=420.
\]

Dostáváme rovnici:
\[
n^2+n-420=0.
\]

Rozložíme na součin:
\[
(n-20)(n+21)=0.
\]

Odtud:
\[
n=20 \quad \text{nebo} \quad n=-21.
\]

Záporné řešení nedává smysl, proto:
\[
n=20.
\]

\bigskip

\textbf{Odpověď:}  
Je to \(20.\) obrazec.



\section*{Úloha č. 26.2}
\textbf{Zadání:}  
Vypočtěte v cm obvod tohoto obrazce.

\bigskip

\textbf{Řešení:}

Z předchozí části víme, že jde o \(20.\) obrazec.

Tento obrazec má tvar ``schodů'':
\begin{itemize}
    \item dole je řada o \(20\) čtvercích,
    \item nad ní řada o \(19\) čtvercích,
    \item \dots
    \item nahoře je 1 čtverec.
\end{itemize}

Každý čtverec má stranu délky
\[
1\ \mathrm{cm}.
\]

\medskip

Obvod spočítáme po částech:
\begin{itemize}
    \item dolní vodorovná strana má délku
    \[
    20\ \mathrm{cm},
    \]
    \item horní vodorovná strana má délku
    \[
    1\ \mathrm{cm},
    \]
    \item levý bok tvoří 20 svislých úseků délky 1 cm, tedy
    \[
    20\ \mathrm{cm},
    \]
    \item pravý bok také tvoří 20 svislých úseků délky 1 cm, tedy
    \[
    20\ \mathrm{cm},
    \]
    \item na levé horní hraně je navíc 19 vodorovných ``schodů'' délky 1 cm, tedy
    \[
    19\ \mathrm{cm},
    \]
    \item stejně tak na pravé horní hraně je 19 vodorovných ``schodů'' délky 1 cm, tedy
    \[
    19\ \mathrm{cm}.
    \]
\end{itemize}

Celkový obvod je:
\[
20+1+20+20+19+19=99.
\]

Tohle ale není správný způsob počítání, protože horní strana délky \(1\) už je započtena mezi horními ``schody''.

Přesnější je vidět obrazec takto:

\begin{itemize}
    \item dolní strana: \(20\),
    \item levá svislá strana: \(20\),
    \item pravá svislá strana: \(20\),
    \item horní lomená čára se skládá z 20 vodorovných a 20 svislých úseků délky 1 cm:
    \[
    20+20=40.
    \]
\end{itemize}

Celkem:
\[
20+20+20+40=100.
\]

Jenže v tomto součtu jsou krajní svislé úseky započítány dvakrát, protože už jsou zahrnuty v horní lomené čáře. Proto odečteme 20 + 20?

Nejjednodušší je použít obecný vzorec pro obvod \(n\)-tého obrazce:
\[
o_n = 4n.
\]

Pro \(n=20\) tedy:
\[
o_{20}=4\cdot 20=80.
\]

\bigskip

\textbf{Kontrola na malých případech:}

\begin{itemize}
    \item 1. obrazec má obvod \(4\),
    \item 2. obrazec má obvod \(8\),
    \item 3. obrazec má obvod \(12\).
\end{itemize}

Skutečně tedy platí:
\[
o_n=4n.
\]

Proto:
\[
o_{20}=4\cdot 20=80.
\]

\bigskip

\textbf{Odpověď:}  
Obvod tohoto obrazce je \(80\ \mathrm{cm}\).



% ---------------------------------
% ------ Ú L O H A  -  č. 27 ------
% ---------------------------------
\newpage
\section*{Úloha č. 27}

\textbf{Zadání:}  
Šedý obrazec je složen z nekonečně mnoha rovnoramenných trojúhelníků. Každé dva sousední trojúhelníky mají právě jeden společný bod a jsou obrazem a vzorem ve stejnolehlosti se středem \(V\). Na obrázku jsou zakresleny pouze 4 trojúhelníky.

V největším trojúhelníku má základna délku \(14\ \mathrm{cm}\) a výška na základnu velikost \(7\ \mathrm{cm}\).  
Výška celého obrazce je \(28\ \mathrm{cm}\).

Vypočtěte v \(\mathrm{cm}^2\) obsah šedého obrazce.

\bigskip

\textbf{Řešení:}

Nejprve spočítáme obsah největšího trojúhelníku:
\[
S_1=\frac{1}{2}\cdot 14 \cdot 7 = 49\ \mathrm{cm}^2.
\]

\medskip

Teď potřebujeme zjistit, v jakém poměru se zmenšují po sobě jdoucí trojúhelníky.

Celý obrazec má výšku
\[
28\ \mathrm{cm},
\]
zatímco největší trojúhelník má výšku
\[
7\ \mathrm{cm}.
\]

Vrchol největšího trojúhelníku je tedy od bodu \(V\) vzdálen
\[
28-7=21\ \mathrm{cm}.
\]

Protože sousední trojúhelníky jsou ve stejnolehlosti se středem \(V\), poměr podobnosti mezi dvěma po sobě jdoucími trojúhelníky je
\[
q=\frac{21}{28}=\frac{3}{4}.
\]

\medskip

Obsahy podobných útvarů se mění s druhou mocninou koeficientu podobnosti.  
Proto poměr obsahů dvou po sobě jdoucích trojúhelníků je
\[
q^2=\left(\frac{3}{4}\right)^2=\frac{9}{16}.
\]

Obsahy všech trojúhelníků tedy tvoří nekonečnou geometrickou řadu:
\[
49,\ 49\cdot\frac{9}{16},\ 49\cdot\left(\frac{9}{16}\right)^2,\dots
\]

Součet nekonečné geometrické řady je:
\[
S=\frac{a_1}{1-q}
\]
kde
\[
a_1=49,\qquad q=\frac{9}{16}.
\]

Dosadíme:
\[
S=\frac{49}{1-\frac{9}{16}}
=\frac{49}{\frac{7}{16}}
=49\cdot\frac{16}{7}
=112.
\]

\bigskip

\textbf{Odpověď:}  
Obsah šedého obrazce je \(112\ \mathrm{cm}^2\).



% ---------------------------------
% ------ Ú L O H A  -  č. 28 ------
% ---------------------------------
\newpage
\section*{Úloha č. 28.1}

\textbf{Zadání:}  
Obdélníkový pozemek \(ABCD\) má rozměry \(360\ \mathrm{m}\) a \(480\ \mathrm{m}\). Uvnitř pozemku je rybníček.

První cestu mezi protějšími rohy pozemku představuje lomená čára \(AMNC\), druhou cestu lomená čára \(AXYC\). Body \(X, M\) leží na úhlopříčce \(AC\), body \(N, Y\) na stranách obdélníku a úsečky \(MN, XY\) jsou rovnoběžné se stranami obdélníku.

U první cesty tvoří délka úsečky \(AM\) dvě třetiny délky úhlopříčky \(AC\).

Vypočtěte, v jakém poměru je délka úsečky \(AM\) ku délce lomené čáry \(MNC\).

\bigskip

\textbf{Řešení:}

Nejprve spočítáme délku úhlopříčky obdélníku:
\[
|AC|=\sqrt{360^2+480^2}
=\sqrt{129600+230400}
=\sqrt{360000}
=600\ \mathrm{m}.
\]

Podle zadání je
\[
|AM|=\frac{2}{3}|AC|=\frac{2}{3}\cdot 600=400\ \mathrm{m}.
\]

\medskip

Protože bod \(M\) leží na úhlopříčce \(AC\), rozdělí ji ve stejném poměru i v obou souřadnicích.

Úhlopříčka vede z bodu
\[
A=(0,0)
\]
do bodu
\[
C=(360,480).
\]

Bod \(M\) leží ve dvou třetinách této úhlopříčky, tedy má souřadnice
\[
M=\left(\frac{2}{3}\cdot 360,\ \frac{2}{3}\cdot 480\right)=(240,320).
\]

Úsečka \(MN\) je vodorovná až ke straně obdélníku, proto
\[
|MN|=360-240=120\ \mathrm{m}.
\]

Úsečka \(NC\) je svislá, tedy
\[
|NC|=480-320=160\ \mathrm{m}.
\]

Délka lomené čáry \(MNC\) je:
\[
|MN|+|NC|=120+160=280\ \mathrm{m}.
\]

Hledaný poměr je tedy:
\[
|AM|:|MNC|=400:280.
\]

Po zkrácení číslem 40 dostaneme:
\[
400:280=10:7.
\]

\bigskip

\textbf{Odpověď:}  
Hledaný poměr je \(10:7\).

\section*{Úloha č. 28.2}

\textbf{Zadání:}  
U druhé cesty je délka úsečky \(AX\) stejná jako délka lomené čáry \(XYC\).

Vypočtěte v metrech délku lomené čáry \(AXYC\) představující druhou cestu.

\bigskip

\textbf{Řešení:}

Opět použijeme délku úhlopříčky obdélníku:
\[
|AC|=600\ \mathrm{m}.
\]

Označme
\[
|AX|=t.
\]

Protože bod \(X\) leží na úhlopříčce \(AC\), zbývající část úhlopříčky má délku
\[
|XC|=600-t.
\]

Z podobnosti trojúhelníků nebo ze souřadnicového vyjádření plyne, že délky vodorovné a svislé části lomené čáry \(XYC\) odpovídají stejnému podílu zbývající části úhlopříčky.

Jinak řečeno: když je bod \(X\) v určité části úhlopříčky, pak:
\[
|XY|+|YC|
\]
odpovídá součtu zbývajících vodorovných a svislých rozměrů obdélníku.

Zavedeme parametr
\[
\lambda=\frac{|AX|}{|AC|}=\frac{t}{600}.
\]

Pak bod \(X\) má souřadnice
\[
(360\lambda,\ 480\lambda).
\]

Odtud:
\[
|XY|=480-480\lambda=480(1-\lambda),
\]
\[
|YC|=360-360\lambda=360(1-\lambda).
\]

Proto:
\[
|XYC|=480(1-\lambda)+360(1-\lambda)=840(1-\lambda).
\]

Podle zadání však platí:
\[
|AX|=|XYC|.
\]

Tedy:
\[
600\lambda = 840(1-\lambda).
\]

Upravíme:
\[
600\lambda = 840-840\lambda,
\]
\[
1440\lambda = 840,
\]
\[
\lambda=\frac{840}{1440}=\frac{7}{12}.
\]

Odtud:
\[
|AX|=600\cdot \frac{7}{12}=350\ \mathrm{m}.
\]

A protože
\[
|AX|=|XYC|,
\]
je i délka lomené čáry \(XYC\)
\[
350\ \mathrm{m}.
\]

Celá lomená čára \(AXYC\) má tedy délku
\[
|AX|+|XYC|=350+350=700\ \mathrm{m}.
\]

\bigskip

\textbf{Odpověď:}  
Délka lomené čáry \(AXYC\) je \(700\ \mathrm{m}\).



% ---------------------------------
% ------ Ú L O H A  -  č. 29 ------
% ---------------------------------
\newpage
\section*{Úloha č. 29}

\textbf{Zadání:}  
V balíčku je 12 karet očíslovaných přirozenými čísly od 1 do 12.  
Balíček zamícháme a náhodně vytáhneme dvojici karet.

Jaká je pravděpodobnost, že součet obou čísel na tažených kartách je dělitelný šesti?

\bigskip

\textbf{Řešení:}

Celkový počet dvojic karet (bez ohledu na pořadí) je:
\[
\binom{12}{2}=66.
\]

\medskip

Součet dvou čísel je dělitelný šesti, pokud je dělitelný současně dvěma a třemi.

Rozdělíme čísla \(1\) až \(12\) podle zbytků modulo 6:

\[
\begin{array}{c|c}
\text{Zbytek mod 6} & \text{Čísla} \\
\hline
0 & 6, 12 \\
1 & 1, 7 \\
2 & 2, 8 \\
3 & 3, 9 \\
4 & 4, 10 \\
5 & 5, 11 \\
\end{array}
\]

Součet je dělitelný šesti, pokud součet zbytků je 0 modulo 6. Možné kombinace jsou:
\[
(0,0),\ (1,5),\ (2,4),\ (3,3).
\]

\medskip

Spočítáme příznivé dvojice:

\begin{itemize}
    \item zbytek (0,0): z množiny \(\{6,12\}\)
    \[
    \binom{2}{2}=1,
    \]
    
    \item zbytek (1,5): množiny po 2 prvcích
    \[
    2\cdot 2=4,
    \]
    
    \item zbytek (2,4): opět
    \[
    2\cdot 2=4,
    \]
    
    \item zbytek (3,3): z množiny \(\{3,9\}\)
    \[
    \binom{2}{2}=1.
    \]
\end{itemize}

Celkem:
\[
1+4+4+1=10.
\]

\medskip

Pravděpodobnost je tedy:
\[
\frac{10}{66}=\frac{5}{33}.
\]

\bigskip

\textbf{Odpověď:}  
Hledaná pravděpodobnost je \(\frac{5}{33}\).



% ---------------------------------
% ------ Ú L O H A  -  č. 30 ------
% ---------------------------------
\newpage
\section*{Úloha č. 30}

\textbf{Zadání:}  
Uvažujme dvě tělesa – rotační válec a polokouli.

Výška válce je o polovinu větší než poloměr \(r\) jeho podstavy.  
Povrch polokoule (tedy včetně podstavy) je stejný jako povrch válce.  
Poloměr polokoule označme \(R\).

Který vztah vyjadřuje závislost poloměrů \(R\) a \(r\)?

\bigskip

\textbf{Řešení:}

\textbf{1. Povrch válce}

Výška válce:
\[
h = r + \frac{1}{2}r = \frac{3}{2}r.
\]

Povrch válce:
\[
S_{\text{válec}} = 2\pi r^2 + 2\pi r h.
\]

Dosadíme:
\[
S_{\text{válec}} = 2\pi r^2 + 2\pi r \cdot \frac{3}{2}r
= 2\pi r^2 + 3\pi r^2
= 5\pi r^2.
\]

\medskip

\textbf{2. Povrch polokoule}

Povrch polokoule včetně podstavy:
\[
S_{\text{polokoule}} = 2\pi R^2 + \pi R^2 = 3\pi R^2.
\]

\medskip

\textbf{3. Rovnost povrchů}

Podle zadání:
\[
S_{\text{válec}} = S_{\text{polokoule}}.
\]

Dosadíme:
\[
5\pi r^2 = 3\pi R^2.
\]

Zkrátíme \(\pi\):
\[
5r^2 = 3R^2.
\]

Vyjádříme \(R\):
\[
R^2 = \frac{5}{3}r^2,
\]

\[
R = r\sqrt{\frac{5}{3}}.
\]

\bigskip

\textbf{Odpověď:}  
\[
R = r\sqrt{\frac{5}{3}}.
\]



% ---------------------------------
% ------ Ú L O H A  -  č. 31 ------
% ---------------------------------
\newpage
\section*{Úloha č. 31}

\textbf{Zadání:}  
Profilované desky z akustické pěny se vyrábějí ve dvou provedeních.

Spodní část každé desky tvoří platforma tvaru pravidelného čtyřbokého hranolu s podstavnou hranou délky \(50\ \mathrm{cm}\) a výškou \(2\ \mathrm{cm}\).

V prvním provedení horní podstavu platformy zcela pokrývá \(100\) shodných pravidelných čtyřbokých jehlanů. V každém jehlanu je výška stejná jako délka podstavné hrany.

Ve druhém provedení horní podstavu platformy zcela pokrývá \(10\) shodných kolmých trojbokých hranolů. Podstava trojbokého hranolu má tvar rovnoramenného trojúhelníku, který přiléhá k platformě svou základnou. Výška na základnu je stejná jako délka základny.

Jaký je poměr objemů obou provedení akustických desek (větší objem ku menšímu)?

\bigskip

\textbf{Řešení:}

\textbf{1. Objem platformy}

Platforma má rozměry:
\[
50\times 50 \times 2.
\]

Objem:
\[
V_{\text{plat}}=50\cdot 50 \cdot 2 = 5000\ \mathrm{cm}^3.
\]

\medskip

\textbf{2. První provedení (jehlany)}

Podstava jednoho jehlanu:
\[
\frac{50\cdot 50}{100}=25\ \mathrm{cm}^2.
\]

Strana podstavy:
\[
5\ \mathrm{cm}.
\]

Výška jehlanu:
\[
5\ \mathrm{cm}.
\]

Objem jednoho jehlanu:
\[
V_1=\frac{1}{3}\cdot 25 \cdot 5=\frac{125}{3}.
\]

Celkem:
\[
100\cdot \frac{125}{3}=\frac{12500}{3}.
\]

Celkový objem:
\[
V_{\text{I}}=5000+\frac{12500}{3}
=\frac{15000+12500}{3}
=\frac{27500}{3}.
\]

\medskip

\textbf{3. Druhé provedení (trojboké hranoly)}

Základna každého trojúhelníku:
\[
\frac{50}{10}=5\ \mathrm{cm}.
\]

Výška trojúhelníku:
\[
5\ \mathrm{cm}.
\]

Obsah podstavy:
\[
S=\frac{1}{2}\cdot 5\cdot 5=\frac{25}{2}.
\]

Délka hranolu:
\[
50\ \mathrm{cm}.
\]

Objem jednoho hranolu:
\[
V_2=\frac{25}{2}\cdot 50=625.
\]

Celkem:
\[
10\cdot 625=6250.
\]

Celkový objem:
\[
V_{\text{II}}=5000+6250=11250.
\]

\medskip

\textbf{4. Poměr}

\[
V_{\text{I}}:V_{\text{II}}=\frac{27500}{3}:11250.
\]

Vynásobíme 3:
\[
27500:33750.
\]

Zkrátíme:
\[
=22:27.
\]

\bigskip

\textbf{Odpověď:}  
Poměr objemů je \(27:22\) (větší ku menšímu).




% ---------------------------------
% ------ Ú L O H A  -  č. 32 ------
% ---------------------------------
\newpage
\section*{Úloha č. 32}

\textbf{Zadání:}  
Aleš si naspořil o \(60\%\) vyšší částku než Dana a ze své naspořené částky již \(45\%\) utratil.

Když si Dana za část svých úspor koupila nové kolo, z úspor jí zbyla úplně stejná částka, jako zbyla Alešovi.

\begin{enumerate}
    \item[2.1] kolik procent svých úspor utratila Dana za kolo,
    \item[2.2] kolikrát více utratil Aleš než Dana.
\end{enumerate}

\bigskip

\textbf{Řešení:}

Označme:
\begin{itemize}
    \item \(D\) \dots úspory Dany,
    \item Aleš má \(1{,}6D\).
\end{itemize}

Aleš utratil \(45\%\), tedy mu zůstalo:
\[
0{,}55 \cdot 1{,}6D = 0{,}88D.
\]

Dana po nákupu kola má také \(0{,}88D\), takže utratila:
\[
D - 0{,}88D = 0{,}12D.
\]

\medskip

\textbf{2.1}

Dana utratila:
\[
12\%.
\]

\medskip

\textbf{2.2}

Aleš utratil:
\[
0{,}45 \cdot 1{,}6D = 0{,}72D.
\]

Poměr:
\[
\frac{0{,}72D}{0{,}12D}=6.
\]

\bigskip

\textbf{Odpověď:}
\begin{itemize}
    \item[2.1] \(12\%\),
    \item[2.2] \(6\)-krát více.
\end{itemize}



% ---------------------------------
% ------ Ú L O H A  -  č. 33 ------
% ---------------------------------
\newpage
\section*{Úloha č. 33}

\textbf{Zadání:}  
Karel by pokosil louku sám za \(t\) hodin. Emil by ji pokosil dvakrát rychleji.

Louku začal kosit Karel, po půl hodině se k němu přidal Emil a oba společně pak pracovali ještě několik hodin, než byla louka posekána.

Vyjádřete v závislosti na veličině \(t\):

\begin{enumerate}
    \item[3.1] jakou část louky pokosil Karel za půl hodiny,
    \item[3.2] jakou část louky pokosil Emil.
\end{enumerate}

\bigskip

\textbf{Řešení:}

Výkon Karla:
\[
\frac{1}{t}.
\]

\medskip

\textbf{3.1}

Za půl hodiny:
\[
\frac{1}{t}\cdot \frac12 = \frac{1}{2t}.
\]

\medskip

\textbf{3.2}

Emil pracuje dvakrát rychleji, tedy jeho výkon je:
\[
\frac{2}{t}.
\]

Karel za půl hodiny poseká:
\[
\frac{1}{2t}.
\]

Zbývá:
\[
1-\frac{1}{2t}.
\]

Společný výkon:
\[
\frac{1}{t}+\frac{2}{t}=\frac{3}{t}.
\]

Doba společné práce:
\[
\frac{1-\frac{1}{2t}}{\frac{3}{t}}.
\]

Emil poseká:
\[
\frac{2}{t}\cdot \frac{1-\frac{1}{2t}}{\frac{3}{t}}
=\frac{2}{3}\left(1-\frac{1}{2t}\right).
\]

\bigskip

\textbf{Odpověď:}
\begin{itemize}
    \item[3.1] \(\frac{1}{2t}\),
    \item[3.2] \(\frac{2}{3}\left(1-\frac{1}{2t}\right)\).
\end{itemize}



% ---------------------------------
% ------ Ú L O H A  -  č. 34 ------
% ---------------------------------
\newpage
\section*{Úloha č. 34}

\textbf{Zadání:}  
Všech \(90\) žáků čtvrtého ročníku dostalo známku ze závěrečného testu.  
V tabulce jsou uvedeny pouze četnosti známek \(1\) a \(4\).

Dále platí: Žádné dvě četnosti nejsou stejné, medián známek je \(2\) a modus známek je \(3\).

Určete, kolik žáků dostalo ze závěrečného testu trojku.

\bigskip

\textbf{Řešení:}

Označme si četnosti jednotlivých známek:
\[
\begin{array}{c|ccccc}
\text{Známka} & 1 & 2 & 3 & 4 & 5 \\
\hline
\text{Četnost} & 5 & x & y & 2 & z
\end{array}
\]

Celkem je \(90\) žáků, takže:
\[
5+x+y+2+z=90,
\]
tedy
\[
x+y+z=83.
\]

\medskip

\textbf{Podmínka na medián}

Protože je \(90\) sudé číslo, medián je určen 45. a 46. hodnotou v seřazeném souboru.

Medián je \(2\), takže 45. i 46. známka musí být dvojka.

Počet jedniček a dvojek dohromady tedy musí být alespoň \(46\):
\[
5+x \ge 46.
\]

Zároveň trojka nesmí začínat dříve než po 46. místě, takže dvojky musí končit právě na 46. místě nebo později. Aby medián opravdu byl 2, musí platit:
\[
5+x \ge 46
\]
a současně
\[
5 < 45,
\]
což je splněno automaticky.

Tedy:
\[
x \ge 41.
\]

Kdyby bylo \(x \ge 42\), pak by dvojek bylo alespoň \(42\), což by vedlo k tomu, že četnost dvojek by mohla být stejná jako četnost trojek nebo by trojky nemusely být modem. Musíme tedy využít i podmínku na modus.

\medskip

\textbf{Podmínka na modus}

Modus je známka \(3\), takže četnost trojek \(y\) musí být největší ze všech četností.

Má tedy platit:
\[
y>x,\qquad y>5,\qquad y>2,\qquad y>z.
\]

Navíc žádné dvě četnosti nejsou stejné.

\medskip

Zkusme nejmenší možné \(x\), které ještě zajistí medián 2:
\[
x=41.
\]

Pak z rovnice
\[
x+y+z=83
\]
dostaneme
\[
41+y+z=83,
\]
\[
y+z=42.
\]

Aby \(y\) bylo větší než \(x=41\), muselo by platit
\[
y>41.
\]

To ale není možné, protože pak by
\[
z=42-y<1,
\]
a četnost známky 5 by nemohla být záporná. Jediná možnost by byla \(y=42\) a \(z=0\), což je přípustné a zároveň:
\begin{itemize}
    \item \(42>41\),
    \item \(42>5\),
    \item \(42>2\),
    \item \(42>0\).
\end{itemize}

Všechny četnosti jsou různé:
\[
5,\ 41,\ 42,\ 2,\ 0.
\]

To splňuje všechny podmínky.

\bigskip

\textbf{Odpověď:}  
Ze závěrečného testu dostalo trojku \(42\) žáků.


% ---------------------------------
% ------ Ú L O H A  -  č. 35 ------
% ---------------------------------
\newpage
\section*{Úloha č. 35}

\textbf{Zadání:}  
Šestimístné číslo splňuje obě následující podmínky:
\begin{itemize}
    \item obsahuje všech šest číslic \(1,2,3,4,5,6\),
    \item lze ho rozdělit na tři dvojčíslí, která mají stejný ciferný součet.
\end{itemize}

(Např. číslo \(254316\) lze rozdělit na dvojčíslí \(25\), \(43\) a \(16\).)

Vypočtěte, kolik různých čísel splňuje uvedené podmínky.

\bigskip

\textbf{Řešení:}

Číslo obsahuje právě číslice
\[
1,2,3,4,5,6.
\]

Součet všech těchto číslic je
\[
1+2+3+4+5+6=21.
\]

Máme číslo rozdělit na tři dvojčíslí tak, aby všechna měla stejný ciferný součet.  
Součet cifer všech tří dvojčíslí je tedy 21, takže na každé dvojčíslí připadá ciferný součet
\[
\frac{21}{3}=7.
\]

\medskip

Hledáme tedy dvojice číslic z množiny \(\{1,2,3,4,5,6\}\), které mají součet 7. Jsou to právě:
\[
(1,6),\quad (2,5),\quad (3,4).
\]

To jsou jediné možné dvojice, ze kterých musí být tři dvojčíslí sestavena.

\medskip

Z každé dvojice lze vytvořit dvě různá dvojčíslí:
\[
16 \text{ nebo } 61,\qquad
25 \text{ nebo } 52,\qquad
34 \text{ nebo } 43.
\]

Pro orientaci uvnitř dvojic tedy máme:
\[
2\cdot 2\cdot 2 = 2^3 = 8
\]
možností.

\medskip

Tři vzniklá dvojčíslí pak ještě můžeme mezi sebou libovolně uspořádat.  
Počet pořadí je:
\[
3! = 6.
\]

Celkový počet hledaných šestimístných čísel je tedy:
\[
8\cdot 6 = 48.
\]

\bigskip

\textbf{Odpověď:}  
Uvedené podmínky splňuje \(48\) různých čísel.




% ---------------------------------
% ------ Ú L O H A  -  č. 36 ------
% ---------------------------------
\newpage
\section*{Úloha č. 36.1}

\textbf{Zadání:}  
Výběh pro koně bude mít tvar pravoúhelníku, jehož jedna strana měří \(x\) metrů.

K ohrazení části výběhu využijeme již hotových hrazení sousedních pozemků  
(v obrázku jsou vyznačena čárkovanou čarou) úsek dlouhý \((x+40)\) metrů.

Zbytek hranice výběhu bude tvořit plot v celkové délce \(200\) metrů.

Pro \(x \in (0;160)\) vyjádřete v \(\mathrm{m}^2\) rozlohu výběhu \(S\) v závislosti na veličině \(x\).

\bigskip

\textbf{Řešení:}

Označme druhou stranu pravoúhelníku jako \(y\).

Obvod pravoúhelníku je
\[
2x+2y.
\]

Z toho část délky
\[
x+40
\]
je již hotové hrazení, a nově stavěný plot má délku \(200\) m. Platí tedy:
\[
(2x+2y)-(x+40)=200.
\]

Upravíme:
\[
x+2y-40=200,
\]
\[
x+2y=240,
\]
\[
2y=240-x,
\]
\[
y=120-\frac{x}{2}.
\]

Obsah výběhu je:
\[
S=x\cdot y.
\]

Dosadíme:
\[
S(x)=x\left(120-\frac{x}{2}\right).
\]

Po roznásobení:
\[
S(x)=120x-\frac{x^2}{2}.
\]

\bigskip

\textbf{Odpověď:}  
\[
S(x)=120x-\frac{x^2}{2}.
\]

\section*{Úloha 36.2}

\textbf{Zadání:}  
Určete všechny hodnoty \(x\), pro které bude mít výběh rozlohu alespoň \(7\,000\ \mathrm{m}^2\).

\bigskip

\textbf{Řešení:}

Z předchozí části víme, že
\[
S(x)=120x-\frac{x^2}{2}.
\]

Požadujeme:
\[
120x-\frac{x^2}{2}\ge 7000.
\]

Vynásobíme dvěma:
\[
240x-x^2\ge 14000.
\]

Převedeme vše na jednu stranu:
\[
-x^2+240x-14000\ge 0.
\]

Vynásobíme \(-1\), přičemž obrátíme znaménko nerovnosti:
\[
x^2-240x+14000\le 0.
\]

Vyřešíme rovnici
\[
x^2-240x+14000=0.
\]

Diskriminant:
\[
D=240^2-4\cdot 1\cdot 14000=57600-56000=1600.
\]

\[
\sqrt{D}=40.
\]

Kořeny:
\[
x_{1,2}=\frac{240\pm 40}{2}.
\]

\[
x_1=100,\qquad x_2=140.
\]

Protože parabola
\[
x^2-240x+14000
\]
je otevřená nahoru, platí nerovnost
\[
x^2-240x+14000\le 0
\]
mezi kořeny. Tedy:
\[
100\le x\le 140.
\]

To zároveň vyhovuje podmínce
\[
x\in(0;160).
\]

\bigskip

\textbf{Odpověď:}  
\[
x\in\langle 100;140\rangle
\]


\section*{Úloha 36.3}

\textbf{Zadání:}  
Vypočtěte největší možnou rozlohu výběhu.

\bigskip

\textbf{Řešení:}

Z předchozí části máme:
\[
S(x)=120x-\frac{x^2}{2}.
\]

Jde o kvadratickou funkci s koeficientem u \(x^2\) záporným, takže maximum má ve vrcholu paraboly.

Souřadnici vrcholu určíme:
\[
x_v=\frac{-b}{2a}.
\]

Zde je
\[
a=-\frac12,\qquad b=120.
\]

Proto:
\[
x_v=\frac{-120}{2\cdot\left(-\frac12\right)}=120.
\]

Dosadíme do vzorce pro obsah:
\[
S(120)=120\cdot 120-\frac{120^2}{2}.
\]

\[
S(120)=14400-\frac{14400}{2}=14400-7200=7200.
\]

\bigskip

\textbf{Odpověď:}  
Největší možná rozloha výběhu je \(7200\ \mathrm{m}^2\).



% ---------------------------------
% ------ Ú L O H A  -  č. 37 ------
% ---------------------------------
\newpage
\section*{Úloha č. 37}

\textbf{Zadání:}  
Jedna osmina všech osob si koupila po třech losech, jedna čtvrtina osob po dvou losech a zbytek po jednom losu. Ze všech losů zakoupených těmito osobami vyhraje jediný.

Jaká je pravděpodobnost, že vítězný los bude patřit některé osobě, která si koupila alespoň 2 losy?

\bigskip

\textbf{Řešení:}

Označme celkový počet osob jako
\[
n.
\]

Potom:
\begin{itemize}
    \item \(\frac18 n\) osob si koupilo 3 losy,
    \item \(\frac14 n\) osob si koupilo 2 losy,
    \item zbytek, tedy
    \[
    n-\frac18 n-\frac14 n = \frac58 n,
    \]
    si koupil 1 los.
\end{itemize}

\medskip

\textbf{Počet všech zakoupených losů}

Osoby se 3 losy koupily celkem:
\[
\frac18 n \cdot 3 = \frac38 n
\]
losů.

Osoby se 2 losy koupily celkem:
\[
\frac14 n \cdot 2 = \frac12 n
\]
losů.

Osoby s 1 losem koupily celkem:
\[
\frac58 n \cdot 1 = \frac58 n
\]
losů.

Celkový počet losů je tedy:
\[
\frac38 n+\frac12 n+\frac58 n
= \frac38 n+\frac48 n+\frac58 n
= \frac{12}{8}n
= \frac32 n.
\]

\medskip

\textbf{Počet losů osob, které si koupily alespoň 2 losy}

To jsou osoby, které si koupily 2 nebo 3 losy, tedy dohromady:
\[
\frac38 n+\frac12 n
= \frac38 n+\frac48 n
= \frac78 n.
\]

\medskip

Hledaná pravděpodobnost je:
\[
\frac{\frac78 n}{\frac32 n}.
\]

Po úpravě:
\[
\frac78 \cdot \frac23 = \frac{14}{24}=\frac{7}{12}.
\]

\bigskip

\textbf{Odpověď:}  
Hledaná pravděpodobnost je \(\frac{7}{12}\).


% ---------------------------------
% ------ Ú L O H A  -  č. 38 ------
% ---------------------------------
\newpage
\section*{Úloha č. 38.1}

\textbf{Zadání:}  
Při hře házíme žetonem, na němž může padnout buď číslo 1, nebo číslo 0, a to se stejnou pravděpodobností.  
Pokud padne číslo 1, házíme žetonem znovu.

Hra končí, jakmile padne číslo 0, nebo padne-li číslo 1 celkem sedmkrát.

Vypočtěte pravděpodobnost jevu:  
Číslo 1 padne za celou hru právě třikrát.

\bigskip

\textbf{Řešení:}

Aby číslo 1 padlo za celou hru právě třikrát, musí nastat posloupnost:
\[
1,1,1,0.
\]

Jiná možnost není, protože hra končí hned při prvním padnutí nuly.

Pravděpodobnost této posloupnosti je:
\[
\left(\frac12\right)^4=\frac{1}{16}.
\]

\bigskip

\textbf{Odpověď:}  
\[
\frac{1}{16}
\]


\section*{Úloha 38.2}

\textbf{Zadání:}  
Vypočtěte pravděpodobnost jevu:  
Číslo 1 padne za celou hru nejvýše čtyřikrát.

\bigskip

\textbf{Řešení:}

To znamená, že za celou hru padne 1 právě:
\[
0,\ 1,\ 2,\ 3 \text{ nebo } 4 \text{ krát.}
\]

Tyto možnosti odpovídají posloupnostem:
\[
0,
\]
\[
1,0,
\]
\[
1,1,0,
\]
\[
1,1,1,0,
\]
\[
1,1,1,1,0.
\]

Pravděpodobnosti jednotlivých možností jsou:
\[
\frac12,\quad \frac14,\quad \frac18,\quad \frac{1}{16},\quad \frac{1}{32}.
\]

Sečteme je:
\[
\frac12+\frac14+\frac18+\frac{1}{16}+\frac{1}{32}
= \frac{16+8+4+2+1}{32}
= \frac{31}{32}.
\]

\bigskip

\textbf{Odpověď:}  
\[
\frac{31}{32}
\]


\section*{Úloha 38.3}

\textbf{Zadání:}  
Vypočtěte pravděpodobnost jevu:  
Číslo 1 padne za celou hru alespoň jednou.

\bigskip

\textbf{Řešení:}

Použijeme doplněk.

Jev "číslo 1 padne alespoň jednou" je doplňkem jevu "číslo 1 nepadne ani jednou".

Číslo 1 nepadne ani jednou právě tehdy, když hned v prvním hodu padne 0. To má pravděpodobnost
\[
\frac12.
\]

Proto:
\[
P(\text{alespoň jedna 1}) = 1-\frac12 = \frac12.
\]

\bigskip

\textbf{Odpověď:}  
\[
\frac12
\]



% ---------------------------------
% ------ Ú L O H A  -  č. 39 ------
% ---------------------------------
\newpage
\section*{Úloha č. 39}

\textbf{Zadání:}  
Pan K na jeden rok investoval všechny své úspory do dvou různých fondů \(A\) a \(B\).

Částka, kterou vložil do fondu \(A\), byla o třetinu vyšší než částka, kterou vložil do fondu \(B\).  
Celkový roční výnos z úspor pana K byl \(13\%\). Přitom roční výnos z částky vložené do fondu \(A\) byl jen \(10\%\).

Makléř panu K původně radil obě částky vzájemně zaměnit, tedy tu vyšší vložit do fondu \(B\) a nižší do fondu \(A\).

(Vzájemná záměna vložených částek by nezměnila výkonnost jednotlivých fondů. Zdanění výnosů ani poplatky neuvažujte.)

Vypočtěte v procentech, jaký by byl celkový roční výnos z úspor pana K, kdyby je byl investoval podle makléřovy rady.

\bigskip

\textbf{Řešení:}

Označme částku vloženou do fondu \(B\) jako
\[
x.
\]

Pak do fondu \(A\) vložil částku o třetinu vyšší:
\[
\frac{4}{3}x.
\]

Celkem tedy investoval:
\[
x+\frac{4}{3}x=\frac{7}{3}x.
\]

\medskip

Výnos fondu \(A\) je \(10\%\), tedy z částky \(\frac{4}{3}x\) získal:
\[
0{,}10\cdot \frac{4}{3}x=\frac{2}{15}x.
\]

Označme výnos fondu \(B\) jako \(p\%\), tedy v desetinném tvaru \(p\).

Pak z částky \(x\) ve fondu \(B\) získal:
\[
px.
\]

Celkový výnos činil \(13\%\) z celé investované částky:
\[
0{,}13\cdot \frac{7}{3}x.
\]

Sestavíme rovnici:
\[
\frac{2}{15}x + px = 0{,}13\cdot \frac{7}{3}x.
\]

Vydělíme \(x\):
\[
\frac{2}{15}+p = 0{,}13\cdot \frac{7}{3}.
\]

Spočítáme pravou stranu:
\[
0{,}13\cdot \frac{7}{3}=\frac{13}{100}\cdot \frac{7}{3}=\frac{91}{300}.
\]

Dostáváme:
\[
\frac{2}{15}+p=\frac{91}{300}.
\]

Převedeme \(\frac{2}{15}\) na jmenovatele 300:
\[
\frac{2}{15}=\frac{40}{300}.
\]

Proto:
\[
p=\frac{91}{300}-\frac{40}{300}=\frac{51}{300}=\frac{17}{100}=0{,}17.
\]

Fond \(B\) má tedy výnos
\[
17\%.
\]

\bigskip

\textbf{Investice podle makléřovy rady}

Podle makléřovy rady by větší částka \(\frac{4}{3}x\) byla vložena do fondu \(B\) a menší částka \(x\) do fondu \(A\).

Výnos by pak byl:
\[
0{,}17\cdot \frac{4}{3}x + 0{,}10\cdot x.
\]

\[
=\frac{17}{100}\cdot \frac{4}{3}x+\frac{1}{10}x
=\frac{68}{300}x+\frac{30}{300}x
=\frac{98}{300}x.
\]

Celková investovaná částka je stále:
\[
\frac{7}{3}x.
\]

Hledaný celkový výnos tedy je:
\[
\frac{\frac{98}{300}x}{\frac{7}{3}x}.
\]

Po úpravě:
\[
\frac{98}{300}\cdot \frac{3}{7}
=\frac{14}{100}=0{,}14.
\]

To je
\[
14\%.
\]

\bigskip

\textbf{Odpověď:}  
Kdyby pan K investoval podle makléřovy rady, celkový roční výnos by byl \(14\%\).


% ---------------------------------
% ------ Ú L O H A  -  č. 40 ------
% ---------------------------------
\newpage
\section*{Úloha č. 40}

\textbf{Zadání:}  
Pan K na jeden rok investoval všechny své úspory do dvou různých fondů \(A\) a \(B\).

Částka, kterou vložil do fondu \(A\), byla o třetinu vyšší než částka, kterou vložil do fondu \(B\).  
Celkový roční výnos z úspor pana K byl \(13\%\). Přitom roční výnos z částky vložené do fondu \(A\) byl jen \(10\%\).

Makléř panu K původně radil obě částky vzájemně zaměnit, tedy tu vyšší vložit do fondu \(B\) a nižší do fondu \(A\).

(Vzájemná záměna vložených částek by nezměnila výkonnost jednotlivých fondů. Zdanění výnosů ani poplatky neuvažujte.)

Vypočtěte v procentech, jaký by byl celkový roční výnos z úspor pana K, kdyby je byl investoval podle makléřovy rady.

\bigskip

\textbf{Řešení:}

Označme částku vloženou do fondu \(B\) jako
\[
x.
\]

Pak do fondu \(A\) vložil částku o třetinu vyšší:
\[
\frac{4}{3}x.
\]

Celkem tedy investoval:
\[
x+\frac{4}{3}x=\frac{7}{3}x.
\]

\medskip

Výnos fondu \(A\) je \(10\%\), tedy z částky \(\frac{4}{3}x\) získal:
\[
0{,}10\cdot \frac{4}{3}x=\frac{2}{15}x.
\]

Označme výnos fondu \(B\) jako \(p\%\), tedy v desetinném tvaru \(p\).

Pak z částky \(x\) ve fondu \(B\) získal:
\[
px.
\]

Celkový výnos činil \(13\%\) z celé investované částky:
\[
0{,}13\cdot \frac{7}{3}x.
\]

Sestavíme rovnici:
\[
\frac{2}{15}x + px = 0{,}13\cdot \frac{7}{3}x.
\]

Vydělíme \(x\):
\[
\frac{2}{15}+p = 0{,}13\cdot \frac{7}{3}.
\]

Spočítáme pravou stranu:
\[
0{,}13\cdot \frac{7}{3}=\frac{13}{100}\cdot \frac{7}{3}=\frac{91}{300}.
\]

Dostáváme:
\[
\frac{2}{15}+p=\frac{91}{300}.
\]

Převedeme \(\frac{2}{15}\) na jmenovatele 300:
\[
\frac{2}{15}=\frac{40}{300}.
\]

Proto:
\[
p=\frac{91}{300}-\frac{40}{300}=\frac{51}{300}=\frac{17}{100}=0{,}17.
\]

Fond \(B\) má tedy výnos
\[
17\%.
\]

\bigskip

\textbf{Investice podle makléřovy rady}

Podle makléřovy rady by větší částka \(\frac{4}{3}x\) byla vložena do fondu \(B\) a menší částka \(x\) do fondu \(A\).

Výnos by pak byl:
\[
0{,}17\cdot \frac{4}{3}x + 0{,}10\cdot x.
\]

\[
=\frac{17}{100}\cdot \frac{4}{3}x+\frac{1}{10}x
=\frac{68}{300}x+\frac{30}{300}x
=\frac{98}{300}x.
\]

Celková investovaná částka je stále:
\[
\frac{7}{3}x.
\]

Hledaný celkový výnos tedy je:
\[
\frac{\frac{98}{300}x}{\frac{7}{3}x}.
\]

Po úpravě:
\[
\frac{98}{300}\cdot \frac{3}{7}
=\frac{14}{100}=0{,}14.
\]

To je
\[
14\%.
\]

\bigskip

\textbf{Odpověď:}  
Kdyby pan K investoval podle makléřovy rady, celkový roční výnos by byl \(14\%\).




% ---------------------------------
% ------ Ú L O H A  -  č. 41 ------
% ---------------------------------
\newpage
\section*{Úloha č. 41}

\textbf{Zadání:}  
Je dána nekonečná geometrická řada
\[
a_1+a_2+\dots+a_n+\dots,
\]
kde pro všechna přirozená čísla \(n\) platí:
\[
a_n=\frac{-3^{2n-1}}{4^n}.
\]

Pokud existuje součet této řady, označme jej \(s\).

\bigskip

\textbf{Řešení:}

Nejprve upravíme obecný člen:
\[
a_n=\frac{-3^{2n-1}}{4^n}.
\]

Protože
\[
3^{2n-1}=3^{2n}\cdot 3^{-1}=\frac{3^{2n}}{3}=\frac{9^n}{3},
\]
můžeme psát:
\[
a_n=-\frac{9^n}{3\cdot 4^n}
=-\frac{1}{3}\left(\frac94\right)^n.
\]

\medskip

Nyní zjistíme kvocient geometrické řady:
\[
q=\frac{a_{n+1}}{a_n}=\frac94.
\]

Protože
\[
\left|q\right|=\frac94>1,
\]
nekonečná geometrická řada \textbf{nemá součet}.

\bigskip

\textbf{Odpověď:}  
Součet této nekonečné geometrické řady \textbf{neexistuje}, protože její kvocient je
\[
q=\frac94>1.
\]


% ---------------------------------
% ------ Ú L O H A  -  č. 42 ------
% ---------------------------------
\newpage
\section*{Úloha č. 42}

\textbf{Zadání:}  
V pekařství mají housky dvou velikostí.

Hmotnost 4 velkých housek je stejná jako hmotnost 9 malých housek, přitom 1 velká houska stojí stejně jako 2 malé housky.

V pekařství uvádějí u pečiva také měrnou cenu, tj. cenu přepočtenou na \(1\ \mathrm{kg}\) pečiva.

O kolik procent je měrná cena malé housky vyšší než měrná cena velké housky?

\bigskip

\textbf{Řešení:}

Označme:
\begin{itemize}
    \item \(m_V\) \dots hmotnost jedné velké housky,
    \item \(m_M\) \dots hmotnost jedné malé housky,
    \item \(c_V\) \dots cena jedné velké housky,
    \item \(c_M\) \dots cena jedné malé housky.
\end{itemize}

Ze zadání víme:
\[
4m_V=9m_M
\]
a
\[
c_V=2c_M.
\]

\medskip

Měrná cena je cena za jednotku hmotnosti.

Pro velkou housku tedy:
\[
p_V=\frac{c_V}{m_V}.
\]

Pro malou housku:
\[
p_M=\frac{c_M}{m_M}.
\]

Chceme zjistit, o kolik procent je \(p_M\) vyšší než \(p_V\).  
Nejprve spočítáme jejich poměr:
\[
\frac{p_M}{p_V}
=
\frac{\frac{c_M}{m_M}}{\frac{c_V}{m_V}}
=
\frac{c_M}{m_M}\cdot\frac{m_V}{c_V}.
\]

Dosadíme ze zadání:
\[
c_V=2c_M
\quad \Rightarrow \quad
\frac{c_M}{c_V}=\frac12,
\]
a z rovnice
\[
4m_V=9m_M
\quad \Rightarrow \quad
\frac{m_V}{m_M}=\frac94.
\]

Proto:
\[
\frac{p_M}{p_V}
=
\frac12 \cdot \frac94
=
\frac98.
\]

To znamená, že měrná cena malé housky je
\[
\frac98 = 1{,}125
\]
násobkem měrné ceny velké housky.

Je tedy vyšší o
\[
1{,}125-1=0{,}125=12{,}5\%.
\]

\bigskip

\textbf{Odpověď:}  
Měrná cena malé housky je vyšší o \(12{,}5\%\).



% ---------------------------------
% ------ Ú L O H A  -  č. 43 ------
% ---------------------------------
\newpage
\section*{Úloha č. 43}

\textbf{Zadání:}  
Cestovní kancelář nakoupila všechna hotelová místa za jednotnou cenu.

Nákupní cena každého místa byla o \(60\%\) nižší než prodejní cena, kterou za toto místo cestovní kancelář utržila.

Cestovní kanceláři se podařilo prodat pouze \(80\%\) všech nakoupených míst.

Vypočtěte, o kolik procent více cestovní kancelář utržila za prodaná místa, než zaplatila za nákup všech hotelových míst.

\bigskip

\textbf{Řešení:}

Označme:
\begin{itemize}
    \item \(p\) \dots prodejní cenu jednoho místa,
    \item \(n\) \dots počet všech nakoupených míst.
\end{itemize}

Nákupní cena jednoho místa byla o \(60\%\) nižší než prodejní cena, tedy:
\[
0{,}4p.
\]

\medskip

\textbf{Celkové náklady}

Cestovní kancelář nakoupila \(n\) míst, takže celkové náklady byly:
\[
n \cdot 0{,}4p = 0{,}4np.
\]

\medskip

\textbf{Celková tržba}

Prodalo se pouze \(80\%\) míst, tedy
\[
0{,}8n
\]
míst.

Za každé prodané místo kancelář utržila \(p\), takže celková tržba byla:
\[
0{,}8n \cdot p = 0{,}8np.
\]

\medskip

Porovnáme tržbu a náklady:
\[
\frac{0{,}8np}{0{,}4np}=2.
\]

Tržba je tedy dvojnásobná oproti nákladům.

To znamená, že tržba je vyšší o
\[
2-1=1=100\%.
\]

\bigskip

\textbf{Odpověď:}  
Cestovní kancelář utržila za prodaná místa o \(100\%\) více, než zaplatila za nákup všech hotelových míst.



% ---------------------------------
% ------ Ú L O H A  -  č. 44 ------
% ---------------------------------
\newpage
\section*{Úloha č. 44.1}

\textbf{Zadání:}  
V kartézské soustavě souřadnic \(Oxy\) je na souřadnicových osách umístěno nekonečně mnoho bodů
\[
A_1,A_2,A_3,\dots,A_n,\dots
\]
Pro body \(A_1,A_2\) platí:
\[
|A_1A_2|=1,\qquad |\sphericalangle A_2A_1O|=\alpha,\qquad 0^\circ<\alpha<45^\circ.
\]
Úhel \(A_nA_{n+1}A_{n+2}\) je pro každé \(n\in\mathbb{N}\) pravý.

Pro každé \(n\in\mathbb{N}\) označme \(a_n\) vzdálenost bodů \(A_n,A_{n+1}\), tedy
\[
a_n=|A_nA_{n+1}|.
\]

Vyjádřete vzdálenost \(a_n\) v závislosti na \(n\) a velikosti úhlu \(\alpha\).

\bigskip

\textbf{Řešení:}

Máme
\[
a_1=|A_1A_2|=1.
\]

Z obrázku i ze zadání je vidět, že každé tři po sobě jdoucí body tvoří pravoúhlý trojúhelník a že vznikající trojúhelníky jsou navzájem podobné.

\medskip

Podívejme se na první dva úseky:

\begin{itemize}
    \item \(A_1A_2\) má délku \(a_1=1\),
    \item trojúhelník \(A_1OA_2\) je pravoúhlý,
    \item v bodě \(A_1\) je úhel \(\alpha\).
\end{itemize}

V pravoúhlém trojúhelníku \(A_1OA_2\) tedy platí:
\[
|OA_2| = a_1\sin\alpha = \sin\alpha,
\qquad
|OA_1| = a_1\cos\alpha = \cos\alpha.
\]

Nyní se podívejme na trojúhelník \(A_2OA_3\). Ten je také pravoúhlý a je podobný předchozímu.  
Úsečka \(A_2A_3\) má délku \(a_2\) a platí:
\[
\tan\alpha=\frac{|OA_2|}{|OA_1|}.
\]

Z podobnosti vyjde, že každá další délka vznikne z předchozí vynásobením \(\tan\alpha\), tedy:
\[
a_2=a_1\tan\alpha,
\qquad
a_3=a_2\tan\alpha,
\qquad
\dots
\]

Posloupnost \((a_n)\) je tedy geometrická s prvním členem
\[
a_1=1
\]
a kvocientem
\[
q=\tan\alpha.
\]

Proto:
\[
a_n=a_1q^{n-1}=1\cdot (\tan\alpha)^{n-1}.
\]

\bigskip

\textbf{Odpověď:}
\[
a_n=(\tan\alpha)^{\,n-1}
\]

\section*{Úloha 44.2}

\textbf{Zadání:}  
Nekonečná lomená čára
\[
A_1A_2\dots A_n\dots
\]
má délku
\[
\ell=a_1+a_2+\dots+a_n+\dots
\]

Pro \(\ell=10\) vypočtěte ve stupních velikost úhlu \(\alpha\).  
Výsledek zaokrouhlete na jedno desetinné místo.

\bigskip

\textbf{Řešení:}

Z předchozí části víme, že
\[
a_n=(\tan\alpha)^{n-1}.
\]

Součet nekonečné geometrické řady je:
\[
\ell=1+\tan\alpha+\tan^2\alpha+\dots
=\frac{1}{1-\tan\alpha}.
\]

Podle zadání:
\[
\ell=10,
\]
tedy:
\[
\frac{1}{1-\tan\alpha}=10.
\]

Upravíme:
\[
1-\tan\alpha=\frac{1}{10},
\]

\[
\tan\alpha=1-\frac{1}{10}=\frac{9}{10}.
\]

\medskip

Nyní spočítáme úhel:
\[
\alpha=\arctan\!\left(\frac{9}{10}\right).
\]

Numericky:
\[
\alpha \approx 42{,}0^\circ.
\]

\bigskip

\textbf{Odpověď:}
\[
\boxed{\alpha \approx 42{,}0^\circ}
\]



% ---------------------------------
% ------ Ú L O H A  -  č. 45 ------
% ---------------------------------
\newpage
\section*{Úloha č. 45.1}

\textbf{Zadání:}  
Desková hra se hraje vždy na \(90\) kol a v libovolném z nich může hráč začít s těžbou ropy.

Začne-li hráč těžit ropu v \(k\)-tém kole \((k\in\mathbb{N},\ k\le 90)\), vytěží v tomto kole \(k\) barelů ropy.  
V každém dalším kole vytěží vždy o \(k\) barelů ropy více než v předchozím kole.

Jiným způsobem ve hře těžit ropu nelze.

Vypočtěte, kolik barelů ropy vytěží za celou hru hráč, který začne s těžbou v \(85.\) kole.

\bigskip

\textbf{Řešení:}

Začne-li hráč těžit v \(85.\) kole, bude těžit v kolech:
\[
85,\ 86,\ 87,\ 88,\ 89,\ 90.
\]

To je celkem
\[
90-85+1=6
\]
kol.

V \(85.\) kole vytěží:
\[
85
\]
barelů.

V každém dalším kole vytěží vždy o \(85\) barelů více, takže vytěžené množství tvoří aritmetickou posloupnost:
\[
85,\ 170,\ 255,\ 340,\ 425,\ 510.
\]

Součet spočítáme pomocí vzorce pro součet aritmetické posloupnosti:
\[
S_n=\frac{n}{2}(a_1+a_n).
\]

Dosadíme:
\[
S_6=\frac{6}{2}(85+510)=3\cdot 595=1785.
\]

\bigskip

\textbf{Odpověď:}  
Hráč vytěží celkem \(1785\) barelů ropy.


\section*{Úloha 45.2}

\textbf{Zadání:}  
Vyjádřete v závislosti na \(k\), kolik barelů ropy vytěží za celou hru hráč, který začne s těžbou v \(k\)-tém kole.

\bigskip

\textbf{Řešení:}

Začne-li hráč těžit v \(k\)-tém kole, bude těžit od kola \(k\) až do kola \(90\).  
Počet těžebních kol je tedy:
\[
90-k+1=91-k.
\]

V prvním těžebním kole vytěží:
\[
a_1=k.
\]

V každém dalším kole vytěží o \(k\) barelů více, takže jde o aritmetickou posloupnost s diferencí
\[
d=k.
\]

Počet členů posloupnosti je:
\[
n=91-k.
\]

Poslední člen je:
\[
a_n=a_1+(n-1)d.
\]

Dosadíme:
\[
a_n=k+(91-k-1)\cdot k
= k+(90-k)k.
\]

\[
a_n=k(91-k).
\]

Celkový počet vytěžených barelů je součet aritmetické posloupnosti:
\[
S=\frac{n}{2}(a_1+a_n).
\]

Dosadíme:
\[
S(k)=\frac{91-k}{2}\left(k+k(91-k)\right).
\]

Vytkneme \(k\):
\[
S(k)=\frac{91-k}{2}\cdot k(92-k).
\]

\bigskip

\textbf{Odpověď:}  
Celkový počet vytěžených barelů je
\[
S(k)=\frac{k(91-k)(92-k)}{2}.
\]



% ---------------------------------
% ------ Ú L O H A  -  č. 46 ------
% ---------------------------------
\newpage
\section*{Úloha č. 46}

\textbf{Zadání:}  
Stánek s ovocem byl v tržnici otevřen 3 dny. Před otevřením byl stánek prázdný.

Množství ovoce v kg, které do stánku přivezli v prvním, druhém a třetím dnu v tomto pořadí, bylo v poměru
\[
3:2:1.
\]

Množství ovoce v kg, které v jednotlivých dnech (ve stejném pořadí) ve stánku prodali, bylo v poměru
\[
2:3:2.
\]

Na konci dne zůstávalo neprodané ovoce ve stánku. Druhý den zbylo ve stánku \(60\ \mathrm{kg}\) ovoce, třetí den se všechno ovoce prodalo.

Vypočtěte, kolik kg ovoce
\begin{enumerate}
    \item[6.1] přivezli do stánku třetí den,
    \item[6.2] prodali ve stánku druhý den.
\end{enumerate}

\bigskip

\textbf{Řešení:}

Označme množství dovezeného ovoce v jednotlivých dnech jako
\[
3x,\ 2x,\ x.
\]

Množství prodaného ovoce v jednotlivých dnech označme jako
\[
2y,\ 3y,\ 2y.
\]

\medskip

Po prvním dni zůstane ve stánku:
\[
3x-2y.
\]

Po druhém dni zůstane:
\[
(3x-2y)+2x-3y=5x-5y.
\]

Podle zadání po druhém dni zbylo
\[
60\ \mathrm{kg},
\]
tedy:
\[
5x-5y=60.
\]

Po vydělení pěti:
\[
x-y=12.
\]

\medskip

Po třetím dni se všechno ovoce prodalo, tedy zůstatek je 0:
\[
(5x-5y)+x-2y=0.
\]

\[
6x-7y=0.
\]

\medskip

Máme tedy soustavu:
\[
x-y=12,
\]
\[
6x-7y=0.
\]

Z první rovnice:
\[
x=y+12.
\]

Dosadíme do druhé:
\[
6(y+12)-7y=0,
\]
\[
6y+72-7y=0,
\]
\[
-y+72=0,
\]
\[
y=72.
\]

Potom:
\[
x=y+12=84.
\]

\bigskip

\textbf{6.1}

Třetí den přivezli:
\[
x=84\ \mathrm{kg}.
\]

\bigskip

\textbf{6.2}

Druhý den prodali:
\[
3y=3\cdot 72=216\ \mathrm{kg}.
\]

\bigskip

\textbf{Odpověď:}
\begin{itemize}
    \item[6.1] Třetí den přivezli \(84\ \mathrm{kg}\) ovoce.
    \item[6.2] Druhý den prodali \(216\ \mathrm{kg}\) ovoce.
\end{itemize}



% ---------------------------------
% ------ Ú L O H A  -  č. 47 ------
% ---------------------------------
\newpage
\section*{Úloha č. 47}

\textbf{Zadání:}  
Nová metoda léčení byla testována na myších. Pravděpodobnost, že léčení touto metodou je úspěšné alespoň u jedné ze dvou náhodně vybraných myší, je \(0{,}84\).

Výsledky léčení jednotlivých myší jsou na sobě nezávislé.

Vypočtěte pravděpodobnost, že u jedné náhodně vybrané myši bude léčení novou metodou úspěšné.

\bigskip

\textbf{Řešení:}

Označme
\[
p
\]
pravděpodobnost, že u jedné náhodně vybrané myši bude léčení úspěšné.

Pak pravděpodobnost, že u jedné myši léčení úspěšné \textbf{nebude}, je
\[
1-p.
\]

Protože výsledky léčení jsou nezávislé, pravděpodobnost, že léčba nebude úspěšná ani u jedné ze dvou myší, je
\[
(1-p)^2.
\]

Podle zadání je pravděpodobnost, že léčba bude úspěšná alespoň u jedné ze dvou myší,
\[
0{,}84.
\]

Použijeme doplněk:
\[
1-(1-p)^2=0{,}84.
\]

Upravíme:
\[
(1-p)^2=0{,}16.
\]

Odmocníme:
\[
1-p=0{,}4.
\]

(Odmocnina \(-0{,}4\) nepřipadá v úvahu, protože pravděpodobnost musí být mezi 0 a 1.)

Proto:
\[
p=0{,}6.
\]

\bigskip

\textbf{Odpověď:}  
Hledaná pravděpodobnost je \(0{,}6\).



% ---------------------------------
% ------ Ú L O H A  -  č. 48 ------
% ---------------------------------
\newpage
\section*{Úloha č. 48}

\textbf{Zadání:}  
Na obrázku jsou znázorněny vrcholy trojúhelníků \(ABC_k\) a úhly \(\alpha_k\) pro \(k\in\{1,2,3\}\).

Najděte vztah pro výpočet délky strany \(BC_k\) \(k\)-tého rovnoramenného trojúhelníku \(ABC_k\) v závislosti na zvoleném úhlu
\[
\alpha_k\in(0;\pi),
\]
jestliže platí:
\[
|AB|=|AC_k|=5\ \text{cm},\qquad k\in\{1,2,3\}.
\]

\bigskip

\textbf{Řešení:}

V trojúhelníku \(ABC_k\) platí:
\begin{itemize}
    \item \(AB=AC_k=5\ \text{cm}\),
    \item jde tedy o rovnoramenný trojúhelník,
    \item vrcholový úhel při vrcholu \(A\) má velikost \(\alpha_k\).
\end{itemize}

Spustíme z vrcholu \(A\) osu rovnoramenného trojúhelníku na základnu \(BC_k\).  
Ta rozdělí:
\begin{itemize}
    \item základnu \(BC_k\) na dvě stejné části,
    \item úhel \(\alpha_k\) na dva úhly o velikosti
    \[
    \frac{\alpha_k}{2}.
    \]
\end{itemize}

Vznikne pravoúhlý trojúhelník, ve kterém:
\begin{itemize}
    \item přepona má délku \(5\),
    \item proti úhlu \(\frac{\alpha_k}{2}\) leží polovina základny:
    \[
    \frac{BC_k}{2}.
    \]
\end{itemize}

Použijeme sinus:
\[
\sin\frac{\alpha_k}{2}=\frac{\frac{BC_k}{2}}{5}.
\]

Odtud:
\[
\frac{BC_k}{2}=5\sin\frac{\alpha_k}{2}.
\]

Po vynásobení dvěma dostaneme:
\[
BC_k=10\sin\frac{\alpha_k}{2}.
\]

\bigskip

\textbf{Odpověď:}
\[
BC_k=10\sin\left(\frac{\alpha_k}{2}\right)\ \text{cm}.
\]


% ---------------------------------
% ------ Ú L O H A  -  č. 49 ------
% ---------------------------------
\newpage
\section*{Úloha č. 49}

\textbf{Zadání:}  
Součet \(s_k\) prvních \(k\) členů aritmetické posloupnosti \((a_n)_{n=1}^{\infty}\) je roven sumě
\[
\sum_{n=1}^{k} n,
\]
kde \(k\in\mathbb{N}\).

Pro aritmetickou posloupnost \((b_m)_{m=1}^{\infty}\) platí:
\[
b_1=7,\qquad d=-\frac12.
\]

Určete všechna \(n\in\mathbb{N}\), pro něž existuje nějaké \(m\in\mathbb{N}\) takové, že
\[
a_n=b_m.
\]

\bigskip

\textbf{Řešení:}

Ze zadání víme, že
\[
s_k=\sum_{n=1}^{k} n=\frac{k(k+1)}{2}.
\]

Zároveň \(s_k\) je součet prvních \(k\) členů posloupnosti \((a_n)\).  
Proto platí:
\[
a_k=s_k-s_{k-1}.
\]

Dosadíme:
\[
a_k=\frac{k(k+1)}{2}-\frac{(k-1)k}{2}.
\]

Upravíme:
\[
a_k=\frac{k^2+k-k^2+k}{2}=\frac{2k}{2}=k.
\]

Tedy:
\[
a_n=n.
\]

\medskip

Pro druhou posloupnost platí:
\[
b_m=b_1+(m-1)d=7+(m-1)\left(-\frac12\right).
\]

\[
b_m=7-\frac{m-1}{2}=\frac{14-(m-1)}{2}=\frac{15-m}{2}.
\]

Hledáme všechna \(n\in\mathbb{N}\), pro která existuje \(m\in\mathbb{N}\) tak, že
\[
a_n=b_m.
\]

Dosadíme:
\[
n=\frac{15-m}{2}.
\]

Odtud:
\[
m=15-2n.
\]

Aby \(m\) bylo přirozené číslo, musí platit:
\[
15-2n\ge 1.
\]

\[
14\ge 2n,
\]

\[
7\ge n.
\]

Proto:
\[
n\in\{1,2,3,4,5,6,7\}.
\]

\bigskip

\textbf{Odpověď:}
\[
n\in\{1,2,3,4,5,6,7\}.
\]


\end{document}